\documentclass[11pt,a4paper]{article}

\usepackage[margin=2.5cm]{geometry}
\usepackage{amsmath,amssymb,amsthm}
\usepackage{booktabs}
\usepackage{longtable}
\usepackage{array}
\usepackage{enumitem}
\usepackage{hyperref}
\usepackage{xcolor}
\usepackage{fancyhdr}
\usepackage{caption}

\hypersetup{colorlinks=true,linkcolor=blue!50!black,urlcolor=blue!50!black,
            citecolor=blue!50!black}

\newtheorem{theorem}{Theorem}[section]
\newtheorem{proposition}[theorem]{Proposition}
\newtheorem{lemma}[theorem]{Lemma}
\newtheorem{corollary}[theorem]{Corollary}
\theoremstyle{definition}
\newtheorem{definition}[theorem]{Definition}
\newtheorem{remark}[theorem]{Remark}
\newtheorem{heuristic}[theorem]{Heuristic}
\newtheorem{observation}[theorem]{Observation}

\newcommand{\Z}{\mathbb{Z}}
\newcommand{\F}{\mathbb{F}}
\newcommand{\Cm}{\mathfrak{Cm}}
\newcommand{\RA}{\mathrm{RA}}
\newcommand{\RRA}{\mathrm{RRA}}
\newcommand{\At}{\mathrm{At}}
\newcommand{\score}{\operatorname{score}}
\newcommand{\sgn}{\operatorname{sign}}
\newcommand{\dom}{\operatorname{Dom}}
\newcommand{\vtwo}{v_2}

\pagestyle{fancy}
\fancyhf{}
\rhead{\small Costas arrays of order 32}
\lhead{\small Consolidated investigation report}
\cfoot{\thepage}

\title{\Large Costas Arrays of Order 32:\\
A Consolidated Report on Algebraic, Relational and Computational Approaches}
\author{Investigation record}
\date{}

\begin{document}
\maketitle

\begin{abstract}
\noindent
Order 32 is, with order 33, the smallest order for which the existence of a
Costas array is unknown. This report consolidates a systematic investigation of
algebraic, relation-algebraic, modular, $p$-adic, matroid-theoretic,
partition-counting and satisfiability-based approaches to the problem, together
with the computational results obtained.

The principal positive results are: (i) a complete classification of the
augmentation ladder from algebraically constructed cores, establishing that no
order-32 Costas array arises as a two-point augmentation of any algebraic
order-30 core; a census-complete classification of the three-point rung, proving that no order-32 Costas array contains an order-29 Costas deletion minor,
with the known score-2 near-solution identified as the unique minimum-defect
object in its class up to the dihedral group; (ii) the observation that this
object factors through a unique derived order-31 Costas array, giving it a clean
genealogy; (iii) a measurement, over the complete order-28 census, showing that
the Costas property is almost never hereditary, and that essentially all
inherited sub-arrays arise from corner deletion; and (iv) a method for
generating the order-28 catalogue from algebraic constructions closed under
corner deletion \emph{and corner addition}, reproducing the published census of
712 arrays in 89 dihedral orbits exactly; and (v) a sound telescoping-sum
propagation rule, proved correct and evaluated under a prospectively frozen,
cost-controlled protocol, which on the complete order-12 six-column benchmark
detects $117{,}880$ infeasible prefixes against $77{,}074$ for the baseline ---
$40{,}806$ additional detections for about $5\%$ extra measured work, with no
rejection of any of the $7{,}632$ completable prefixes.

The principal negative results are a sequence of measured refutations of
proposed structural attacks. Each is recorded with the quantitative evidence
that closes it, so that the same ground need not be covered again. A ranked
programme of future work is given in Section~\ref{sec:future}.
\end{abstract}

\tableofcontents
\newpage

%=====================================================================
\section{Scope and how to read this report}
\label{sec:scope}

This document records an investigation into the existence question for Costas
arrays of order 32. It is organised so that the two kinds of result are kept
separate and both are auditable.

\textbf{Positive results} (Sections~\ref{sec:ladder}--\ref{sec:heredity}) are
finite computations whose conclusions are theorems modulo the correctness of the
software, and every one is accompanied by an independent verification path.

\textbf{Negative results} (Sections~\ref{sec:structural}--\ref{sec:partition})
are refutations of proposed methods. These occupy more space than the positive
results, deliberately. A large part of the value of this investigation lies in
knowing which attractive-looking approaches do not work and precisely why, since
each was proposed in good faith and each consumed effort before being closed.
Every closure is quantitative: a measurement, a counterexample, or a proof of
domination by an existing method.

Throughout, a claim is marked as \emph{measured} when it rests on a computation
reported here, \emph{cited} when it rests on the literature, and
\emph{conjectural} when it is neither.

%=====================================================================
\section{Notation}
\label{sec:notation}

\begin{longtable}{@{}l p{0.78\textwidth}@{}}
\toprule
Symbol & Meaning \\
\midrule
\endhead
$n$ & Order of a Costas array; the board is $n\times n$. Usually $n=32$. \\
$h$ & A permutation $h:\{0,\dots,n-1\}\to\{0,\dots,n-1\}$, written zero-based.
      $h_i$ is the \emph{height} (row) of the dot in column $i$. \\
$\Delta_{m,i}$ & $h_{i+m}-h_i$, the vertical difference at horizontal
      separation $m$ starting at column $i$. \\
$m$, $d$ & Horizontal separation (distance) between two columns. $d$ is used in
      solver contexts, $m$ in partition-counting contexts; they denote the same
      quantity. \\
$e$ & A signed vertical difference; $e\in\mathcal V$. \\
$g$ & An absolute vertical distance, $g=|e|$. \\
$\mathcal V$ & The signed palette $\{-(n-1),\dots,-1,1,\dots,n-1\}$,
      $|\mathcal V| = 2n-2$; for $n=32$, $|\mathcal V|=62$. \\
$L_m$ & $n-m$, the number of column pairs at separation $m$. \\
$\score(h)$ & Costas defect: the number of repeated $(d,\Delta)$ pairs,
      $\sum_{d,e}\max(0,\chi_{d,e}-1)$. $\score(h)=0$ iff $h$ is Costas. \\
$\chi_{m,e}$ & $\#\{i: \Delta_{m,i}=e\}$, the occupancy of signed value $e$
      at separation $m$. \\
$y_{m,g}$ & $\chi_{m,g}+\chi_{m,-g}\in\{0,1,2\}$, magnitude occupancy. \\
$T_m$ & $\sum_i \Delta_{m,i}$, the telescoping row sum at separation $m$. \\
$N_q(r)$ & $\#\{e\in\mathcal V : e\equiv r \pmod q\}$, the signed residue
      capacity. \\
$\vtwo(x)$ & The 2-adic valuation of $x$; $\vtwo(0)=\infty$. \\
$D_i$ & Domain of the height variable $h_i$ during propagation. \\
$E_{i,j}$, $D_{ij}$ & Domain of the edge (difference) variable $c_{ij}=h_j-h_i$. \\
$k$ & Number of columns and rows deleted in an augmentation rung; the core has
      order $n-k$. \\
$K$ & A Costas core array of order $n-k$. \\
$C$, $R$ & The deleted column set and deleted row set, each of size $k$. \\
$\delta,\epsilon$ & Compressed displacement of a core pair: column distance and
      row difference \emph{within the core}. \\
$D_4$ & The dihedral group of order 8 acting on the square; the Costas
      symmetry group. \\
$W_1(p,g,c)$ & Welch construction, order $p-1$: $h_i = g^{i+c}\bmod p - 1$. \\
$W_2,W_3$ & $W_1$ with one or two corner dots deleted, orders $p-2$, $p-3$. \\
$G_2(q,\alpha,\beta)$ & Golomb construction over $\F_q$, order $q-2$:
      dot at $(i,j)$ iff $\alpha^i+\beta^j=1$. \\
$G_3,G_4,G_5$ & $G_2$ with corner dots deleted, orders $q-3,q-4,q-5$. \\
$\Cm(G)$ & The complex algebra of a group $G$: universe $\wp(G)$, composition
      the sumset, converse negation, identity $\{0\}$. \\
$\At(\mathfrak A)$ & The atom structure of an atomic relation algebra. \\
$\RA_n$ & Maddux's class of relation algebras with an $n$-dimensional
      relational basis. \\
$\RRA$ & The class of representable relation algebras. \\
$C_{31}$ & The unique order-31 Costas array reachable by one-dot augmentation
      of the algebraic order-30 cores (one $D_4$ orbit). \\
$N_{32}$ & The known order-32 permutation of defect 2; the corner extension of
      $C_{31}$. \\
\bottomrule
\end{longtable}

%=====================================================================
\section{Glossary}
\label{sec:glossary}

\begin{description}[style=nextline,leftmargin=0pt]

\item[Costas array] A permutation matrix of order $n$ such that all
$\binom{n}{2}$ vectors connecting pairs of dots are distinct. Equivalently, for
every horizontal separation $m$, the values $\Delta_{m,i}$ are pairwise
distinct.

\item[Defect / score] The number of repeated displacement vectors. A Costas
array has defect 0. The order-32 object $N_{32}$ has defect 2.

\item[Distinct difference configuration (DDC)] A set of points in $\Z^2$ whose
connecting vectors are all distinct. Costas arrays are the permutation subclass;
sonar sequences and Golomb rectangles are related bounded-shape variants.

\item[Rainbow / anti-Ramsey] A colouring in which every colour is used at most
once on the edge set in question. The Costas condition is a rainbow condition on
each horizontal-separation class, not a monochromatic-avoidance condition. This
distinguishes it from the Ramsey regime where relation-algebraic methods succeed.

\item[Sporadic array] Under the taxonomy used here, a Costas array not
accounted for by a recognised generated or emergent construction. The category is
epistemic rather than intrinsic: it shrinks whenever a new construction is found.
Order 28 is the first order above 5 with no sporadic arrays. At order 29, 16
arrays in two $RW_0$ classes lie outside the Golomb and Welch algebraic families
but are \emph{emergent}, not sporadic, having been produced by the semi-empirical
Rickard method; order 29 likewise has no sporadic arrays.

\item[Augmentation rung] The finite search over all order-32 permutations formed
by placing a Costas core of order $32-k$ in a compressed sub-board and adding
$k$ dots. Parameterised by $k$.

\item[Corner deletion / addition] Removing (adding) a dot at a corner of the
board. This is the only deletion that preserves the Costas property in general,
because the resulting compression is a translation rather than a reordering.

\item[Core-first pruning] In a rung search, testing the placed core's own
displacements before considering the added dots. Exact for defect 0, since a
repeated displacement among a subset is a repeated displacement of the whole.

\item[Vertex potential / cocycle formulation] Representing the array by heights
$h_i$ and edge colours $c_{ij}=h_j-h_i$, so that $c_{ik}=c_{ij}+c_{jk}$. This
reduces $\binom{n}{2}$ edge variables to $n$ vertex variables.

\item[Relational basis] Maddux's finite-dimensional approximation to
representability: a set of atomic networks on $n$ nodes closed under a
one-node amalgamation condition. Membership in $\RA_n$.

\item[Atomic network] A labelling of the ordered pairs of a node set by atoms,
satisfying identity, converse and the triangle condition. The combinatorial
shadow of a representation.

\item[Board-relative network] An atomic network on a \emph{named, ordered} node
set in which the horizontal component of each label is forced to equal the
column gap and the vertical component is bounded. This is what a Costas array
actually is; ordinary relation-algebra networks do not capture the board.

\item[Cube-and-conquer] A SAT paradigm: split the formula into many cubes with a
lookahead solver, solve each with CDCL, and compose the per-cube proofs. The
method behind the largest certified mathematical computations.

\item[DRAT / LRAT] Machine-checkable proof formats for SAT unsatisfiability
claims. Absent from a solver, an UNSAT verdict is a computational observation
rather than a certificate.

\end{description}

%=====================================================================
\section{The problem and the known landscape}
\label{sec:landscape}

\subsection{Status of the existence question}

Costas arrays are known to exist for all orders $n \le 31$. Orders 32 and 33 are
the smallest for which none is known. Exhaustive enumeration has been carried
out for all orders up to 29; orders 30 and 31 have never been exhaustively
enumerated, although arrays are known at both.

Reported enumeration costs (cited) are: order 26, $29.52$ single-processor
years; order 27, approximately 30 CPU-years across 120 machines over three
months; order 28, about 70 CPU-years; order 29, $366.55$ CPU-years. The observed
growth of the search is slightly worse than a factor of 5 per increment of $n$.

\subsection{Why orders 32 and 33 specifically}

The algebraic constructions produce arrays of orders $p-1$, $p-2$, $p-3$ for $p$
prime and $q-2$, $q-3$, and conditionally $q-4$ and $q-5$, for $q$ a prime
power. The orders reachable near 32 are therefore determined by the prime powers
in the window:

\begin{center}
\begin{tabular}{@{}lccccccccc@{}}
\toprule
order & 28 & 29 & 30 & \textbf{31} & \textbf{32} & \textbf{33} & 34 & 35 & 36 \\
\midrule
direct construction & $29{-}1$ & $31{-}2$ & $31{-}1$ & --- & --- & --- &
  $37{-}3$ & $37{-}2$ & $37{-}1$ \\
dots from nearest seed & 0 & 0 & 0 & 1 & 2 & 2 & 0 & 0 & 0 \\
\bottomrule
\end{tabular}
\end{center}

Since $33 = 3\cdot 11$, $34 = 2\cdot 17$, $35 = 5\cdot 7$ and $36 = 6^2$ are all
non-prime-powers, the only prime power in the window above 31 is 37, which
reaches down to 34 unconditionally and to 33 and 32 only through the conditional
$q-4$ and $q-5$ constructions, whose conditions fail. Orders 31, 32 and 33 are
therefore the only orders in $[28,36]$ with no direct construction, at distances
1, 2 and 2 from the nearest seed. Order 31 succeeds at distance 1; orders 32 and
33 do not at distance 2.

\begin{proposition}[Automatic loss-free parameterisation]
\label{prop:polydual}
For the standard Welch and Golomb parameterisations to induce an order-$N$
permutation without deleting field elements or selecting a proper output window,
one requires
\[
q-1=N \qquad\text{or}\qquad q-2=N,
\]
respectively. Equivalently, these are precisely the cases in which the standard
parameter domain has the required cardinality automatically.

For $N=32$ the required field orders are therefore $q=33$ and $q=34$. Neither is
a prime power, so neither finite field exists. Consequently there is no
automatically loss-free Welch or Golomb parameterisation at order 32.
\end{proposition}

The next available field order is $q=37$. A construction over this field must
therefore impose additional structure restricting its values to an appropriate
width-32 window. Under the heuristic model in which the 32 outputs are
independently uniform over the 36 available nonzero field values, the probability
of filling one prescribed window injectively is
\[
P_{\mathrm{window}} \;=\; \frac{32!}{36^{32}} \;=\; 4.1542\times10^{-15}.
\]
This factor is \emph{not} a non-existence theorem for structured polynomial
families. It quantifies why uninformed sampling over the next available field is
ineffective.

The automatic loss-free condition is the polynomial dual of the prime-power gap
visible in the construction table above. The associated window factor is a
separate random-output heuristic explaining why unstructured sampling beyond that
gap is ineffective. Section~\ref{sec:coords} records the computation.

\begin{remark}
Nothing about $32 = 2^5$ plays any role. The arithmetic that distinguishes order
32 is the distribution of prime powers in $\{33,\dots,37\}$. This is worth
stating explicitly because the binary structure of 32 motivated several of the
approaches refuted in Section~\ref{sec:structural}, and in every case the
binary structure turned out to be irrelevant.
\end{remark}

\subsection{The constraint any non-existence proof must satisfy}

Costas arrays exist at orders 31 and 36. If none exists at 32, the counting
function is non-monotonic in a way tied to prime-power gaps. Consequently:

\begin{observation}[Order-sensitivity requirement]
\label{obs:uniformity}
A successful non-existence argument must contain some feature capable of
distinguishing order 32 from nearby orders at which Costas arrays exist, notably
31 and 36. Bounds or relaxations whose hypotheses and numerical strength vary
smoothly and monotonically across this range are unlikely to suffice unless they
include an arithmetic, boundary, or discrete threshold effect.
\end{observation}

This is a research heuristic, not an impossibility theorem. A theorem may be
uniform in $n$ while its numerical or arithmetic conclusion varies with $n$: a
uniform bound $C(n)\le F(n)$ could in principle satisfy $F(31)\ge1$, $F(32)<1$,
$F(36)\ge1$, and Bruck--Ryser is exactly a uniform theorem whose arithmetic
hypothesis distinguishes orders. The existence of arrays at 31 and 36 therefore
does not rule out a general theorem implying non-existence at 32. What the
measurements below establish is that each \emph{particular} relaxation tested
returns the same verdict at 31, 32 and 36.

This observation provides a diagnostic for the particular relaxations examined
in Section~\ref{sec:structural}, not an impossibility theorem. In each case
tested, the relaxation lacked sufficient order-sensitive strength and returned
essentially the same verdict at orders 31, 32 and 36. A stronger density theorem
or arithmetic bound capable of distinguishing these orders remains possible.

\subsection{The three-branch decomposition}

Any proof of non-existence must dispose of three classes.

\begin{enumerate}[label=(\roman*)]
\item \textbf{Periodic arrays.} Golomb and Moreno proved in 1996 that an order-$n$ circular Costas sequence exists only when $n+1$ is prime. Since 33 is composite,
\emph{no circular Costas array of order 32 exists}; this branch is closed. The full Golomb--Moreno conjecture, that every circular Costas array is Welch, was proved in 2015 by
Muratovi\'c-Ribi\'c, Pott, Thomson and Wang, in its equivalent form that every Costas polynomial over a
prime field is a monomial; since circularity forces $n+1$ to be prime, this covers every circular
sequence. The remaining open statement is the Golomb--Taylor conjecture of
1984, that a Costas sequence is singly periodic if and only if it is Welch;
Moreno proved in 1992 that the stronger \emph{shifting property} characterises
the Welch construction, leaving the gap between single periodicity and the
shifting property open.

\item \textbf{Constructed and augmented arrays.} This is where the present
investigation makes progress; see Section~\ref{sec:ladder}.

\item \textbf{Sporadic arrays.} Everything else. This branch is defined by
absence of structure and is therefore immune to algebraic attack by
construction. Section~\ref{sec:heredity} shows that branch (ii) is thinner than
expected, which makes branch (iii) larger, not smaller.
\end{enumerate}

\begin{remark}
The decomposition is a partition only because branch (iii) is the complement.
The category ``sporadic'' is epistemic. During this investigation a
catalogue-generation bug produced an array that appeared unreachable by
construction until an additional closure operation was added; nothing about the
array changed. The same caution applies to published sporadic counts, although
those rest on the enumeration authors' own classification and not on any
generator built here.
\end{remark}

%=====================================================================
\section{Computational baseline}
\label{sec:baseline}

\subsection{Exact solver engineering}

The starting point was a depth-first search on vertex potentials that allocated
per-node vectors and sorted candidate values at every node, achieving
$5.84\times 10^5$ nodes/s at order 32.

For $n\le 32$ the set of used differences at a fixed horizontal distance lies in
$[-(n-1),n-1]$, that is $2n-1 \le 63$ values, and therefore fits in a single
64-bit word. Rewriting with one \texttt{uint64\_t} mask per distance, no
per-node allocation and no sort gives the same search tree and the following
measured speedup:

\begin{center}
\begin{tabular}{@{}lrr@{}}
\toprule
task & original & bitmask rewrite \\
\midrule
$n=13$ exhaustive (18{,}290{,}044 nodes, identical) & $8.55$ s & $1.26$ s \\
$n=32$, 20 s pilot & $5.84\times10^5$ n/s & $4.34\times10^6$ n/s \\
\bottomrule
\end{tabular}
\end{center}

\subsection{Growth law}

Exhaustive enumeration with the vertical-complement symmetry $p_0\le
\lfloor(n-1)/2\rfloor$ gives the following measured tree sizes. Solution counts
were validated against the published census in every case.

\begin{center}
\begin{tabular}{@{}rrrr@{}}
\toprule
$n$ & nodes & ratio & time (1 core) \\
\midrule
11 & 789{,}512 & --- & 0.04 s\\
12 & 3{,}488{,}681 & 4.42 & 0.21 s\\
13 & 18{,}290{,}044 & 5.24 & 1.26 s\\
14 & 83{,}176{,}085 & 4.55 & 19.6 s\\
15 & 435{,}090{,}281 & 5.23 & 82.9 s\\
16 & 2{,}001{,}473{,}413 & 4.60 & 237.9 s\\
\bottomrule
\end{tabular}
\end{center}

The geometric mean ratio is $\mathbf{4.80}$ per order, consistent with the
published observation of slightly worse than $5\times$ per order. Extrapolating
from the measured order-29 cost of $366.55$ CPU-years gives roughly
$5\times10^4$ CPU-years on 2010-era hardware for order 32, or of order
$10^3$--$10^4$ modern core-years.

\subsection{Order 17 as a calibration point}

Order 17 exhaustive enumeration was measured by splitting on the nine
symmetry-reduced top-level branches $p_0 = 0,\dots,8$:

\begin{center}
\begin{tabular}{@{}lrrr@{}}
\toprule
branch & nodes & time & rate \\
\midrule
$p_0=0$ & $1.099\times10^9$ & 95.8 s & $1.15\times10^7$/s \\
$p_0=1$ & $1.138\times10^9$ & 94.8 s & $1.20\times10^7$/s \\
$p_0=2$ & $1.144\times10^9$ & 97.0 s & $1.18\times10^7$/s \\
\bottomrule
\end{tabular}
\end{center}

Nine branches at that rate is approximately $10^{10}$ nodes and
\textbf{about 14 minutes on one core}. This figure is used in
Section~\ref{sec:sat} to calibrate distributed SAT approaches.

\subsection{Fail-closed discipline}\label{sec:discipline}

Two defects in the original solver are recorded because they are easy to
reproduce in any successor. First, it had no completeness flag, so a
time-limited subtree was indistinguishable from an exhausted one --- fatal for
any non-existence claim. Second, the vertical-complement symmetry was silently
skipped whenever a prefix was supplied, since the guard tested $k=0$ and $k$
never reaches 0 under a prefix; a prefix manifest must therefore be
symmetry-reduced itself. Both are corrected in the released solver, which emits
one of \texttt{EXHAUSTED}, \texttt{WITNESS\_FOUND}, \texttt{TIME\_LIMIT},
\texttt{NODE\_LIMIT}, \texttt{INVALID\_PREFIX} or \texttt{INTERNAL\_ERROR}, and
whose merge tool refuses an exhaustive verdict unless every unit reports
\texttt{EXHAUSTED}, all units agree on parameters and catalogue hash, and the
coverage is exactly $0..N{-}1$ with no gap or duplicate.

%=====================================================================
\section{The augmentation ladder}
\label{sec:ladder}

\subsection{Formulation}

\begin{definition}[Augmentation rung]
Fix $k$ and a Costas core $K$ of order $n-k$. Choose a set $C$ of $k$ columns
and a set $R$ of $k$ rows of the $n\times n$ board. Place the core dots at the
surviving columns and rows in increasing order, and place $k$ further dots as a
bijection $C\to R$. The rung is the exhaustive search over all
$(K,C,R,\mu)$.
\end{definition}

Two remarks on the formulation. First, when $C$ and $R$ both consist of
end-blocks the compression is a translation and the core's displacements are
inherited unchanged; this \emph{contiguous} sub-case is the actual construction
mechanism and is reported separately. Second, core-first pruning --- rejecting a
placement as soon as the placed core carries a repeated displacement --- is
exact for defect 0, because the core dots are a subset of the $n$ dots. It is
\emph{not} exact for defects above 0, so the reported histograms above score 0
count only core-intact configurations.

\subsection{Catalogue generation and census verification}
\label{sec:catalogue}

The algebraic possibilities at each order near 32 are determined by the prime
powers available. For order 30 they are $W_1(31)$ and $G_2(32)$; no other family
applies, since $W_2$ needs $p=32$, $W_3$ needs $p=33$, $G_3$ needs $q=33$, $G_4$
needs $q=34$ and $G_5$ needs $q=35$, none of which is prime or a prime power.

A first generator seeded $W_1(29)$, $W_1(31)$, $G_2(31)$ and $G_2(32)$ and
closed \emph{downward} under corner deletion. It produced 704 order-28 arrays in
88 dihedral orbits against a published census of 712 in 89, and 660 order-30
arrays.

\begin{observation}[The missing orbit]
\label{obs:missing}
The deficit was exactly one dihedral orbit and was caused by the closure being
downward-only. $G_2(29)$ has order $q-2 = 27$ and can reach order 28 only
\emph{upwards}. Adding $G_2(29)$ as a seed and corner \emph{addition} to the
closure recovers the missing orbit, whose canonical representative is
\[
(0,16,21,12,18,22,23,26,10,7,24,14,27,6,17,5,25,11,20,2,4,19,15,13,8,1,9,3).
\]
Eight configurations land in this orbit, including $G_2(29,\alpha=3,\beta=3)$
with a corner dot at $(0,27)$ and $G_2(29,\alpha=3,\beta=10)$ with one at
$(0,0)$.
\end{observation}

With the corrected closure the generator produces:

\begin{center}
\begin{tabular}{@{}lrrl@{}}
\toprule
order & generated & published census & status \\
\midrule
27 & 192 / 27 orbits & 204 / 29 & 2 orbits short \\
\textbf{28} & \textbf{712 / 89 orbits} & \textbf{712 / 89} & \textbf{complete} \\
29 & 148 / 21 orbits & 164 / 23 & 16 arrays short \\
30 & 664 / 85 orbits & never enumerated & --- \\
\bottomrule
\end{tabular}
\end{center}

\begin{proposition}[Completeness of the order-28 catalogue]
The generated order-28 catalogue is the complete census.
\end{proposition}
\begin{proof}
The published enumeration is exhaustive with 712 arrays. The generated
catalogue contains 712 distinct permutations of $\{0,\dots,27\}$, each verified
to have Costas defect 0, and is closed under $D_4$ with 89 orbits. Two sets of
equal finite cardinality, one contained in the other, are equal.
\end{proof}

\begin{remark}
Order 29 landing on exactly $164-16=148$ is an internal consistency check: the
literature records precisely 16 order-29 arrays not accounted for by the
Golomb and Welch constructions, and the corrected generator finds every array
that is. Order 27 falling two orbits short is consistent with the $T_4(31,1)$
family and the single known sporadic array, neither of which the closure
reaches.
\end{remark}

\subsection{The rungs}
\label{sec:rungs}

All rungs were run with core-first pruning, independent re-scoring of every
reported low-defect configuration, and a fail-closed coverage manifest.

\begin{center}
\begin{tabular}{@{}rlrrrr@{}}
\toprule
$k$ & core order & cores & measured placements & Costas found & minimum defect \\
\midrule
1 & 30 (algebraic) & 664 & $6.3\times10^5$ & \textbf{8} & 0 \\
2 & 30 (algebraic) & 664 & $2.091\times10^{7}$ & \textbf{0} & \textbf{2} \\
3 & 29 (\emph{census-complete}) & 164 arrays & $5.658\times10^{8}$ & \textbf{0} & 4 \\
4 & 28 (\emph{census-complete}) & 712 arrays & $1.151\times10^{11}$ & \textbf{0} & --- \\
\bottomrule
\end{tabular}
\end{center}

\paragraph{Rung $k=1$ (order 30 $\to$ 31).}
Exactly 8 order-31 Costas arrays are found, forming exactly one dihedral orbit.
Hence there is, up to the square symmetries, \emph{exactly one} order-31 Costas
array reachable by one-dot augmentation of the algebraic order-30 cores. Defect
histogram: $0{:}8$, $1{:}0$, $2{:}24$, $3{:}16$, $4{:}124$, $5{:}112$,
$6{:}596$.

\paragraph{Rung $k=2$ (order 30 $\to$ 32).}
The engine processed $20{,}911{,}360=85\binom{32}{2}^2$ core placements in 33 seconds on one core: one representative of each of the catalogue's 85
dihedral orbits, in every choice of two inserted columns and two inserted rows. Each placement admits $2!$ bijections of the
added dots, and by $D_4$-equivariance (the remark on canonical mode below) the search covers all $664\binom{32}{2}^2\,2!=326{,}709{,}248$ a priori configurations. It
finds \textbf{zero} Costas arrays. (Earlier statements of this count gave $324{,}741{,}120$, the value for a superseded 660-array catalogue.)
Defect histogram: $0{:}0$, $1{:}0$, $2{:}8$, $3{:}0$, $4{:}24$, $5{:}16$,
$6{:}40$, $7{:}56$, $8{:}104$, $9{:}168$, $10{:}584$.

The eight defect-2 configurations were dumped and re-verified with an
independent scorer: all are valid permutations, all have exact defect 2, and
they are \emph{precisely} the dihedral orbit of the known candidate $N_{32}$.
Nothing else in the space attains 2, and defects 1 and 3 are empty. The
contiguous sub-case comprises 11{,}880 configurations with the same optimum.

\paragraph{Rung $k=3$ (order 29 $\to$ 32): census-complete.}
Reported in full below, after the $k=4$ rung, since it postdates the
reconstruction of the order-29 census.

\paragraph{Rung $k=4$ (order 28 $\to$ 32): completed.}
Run in canonical mode as 89 deterministic work units, one dihedral orbit each,
against the census-complete catalogue (\texttt{CENSUS ... MATCH}, catalogue hash
\texttt{b64cc86a0058f1b3}). Independently merged with the fail-closed tool:

\begin{center}
\begin{tabular}{@{}lr@{}}
\toprule
merge verdict & \texttt{EXHAUSTED\_AND\_VERIFIED} \\
units & 89 of 89, all \texttt{EXHAUSTED} \\
core placements & $115{,}087{,}822{,}400$ \\
survivors of the core-first test & $2{,}401$ \\
configurations evaluated & $57{,}624$ \\
\textbf{Costas arrays found} & \textbf{0} \\
aggregate CPU & $2{,}389.8$ s; longest unit $41.1$ s \\
\bottomrule
\end{tabular}
\end{center}

\noindent The engineering rate was $3.79\times10^7$ core placements per second,
achieved by bucketing the core-pair test by column displacement and replacing the
per-placement clear with a generation stamp. Because order 28 has no sporadic
arrays, ``census-complete'' and ``all algebraic'' coincide here, so this is the
first rung whose negative statement quantifies over \emph{every} core of its
order rather than an algebraic subset.

\begin{remark}[Column convention, corrected]
\label{rem:colconv}
Earlier drafts of this table mixed two quantities under one heading: the $k=4$
row reported \emph{measured} core placements from the run, while the $k\le3$ rows
reported an a priori configuration-space count. Re-running $k=2$ and $k=3$ from a
single binary exposed the discrepancy --- $k=2$ measures $2.09\times10^{7}$
placements against a recorded $3.25\times10^{8}$, and $k=3$ measures
$5.17\times10^{8}$ against a recorded $1.77\times10^{10}$. The column now records
measured placements throughout. The scientific conclusions are unaffected: both
reruns reproduce the recorded outcomes exactly (zero Costas arrays; minimum
defect 2 attained uniquely at $k=2$, minimum defect 4 at $k=3$).
\end{remark}

\paragraph{Rung $k=3$ is now census-complete (items F2 and F1b, both closed).}
The complete order-29 census was reconstructed from Table 2 of the enumeration
paper, which lists the 23 lexicographically minimal equivalence-class
representatives with their construction method. Every representative was
independently validated as an order-29 Costas permutation and expanded under
$D_4$, reproducing all published figures exactly: \textbf{164 arrays, 23
equivalence classes, 5 symmetric}, with per-method counts
$G_0\,4$, $G_2\,64$, $G_3\,24$, $W_0\,24$, $W_2\,32$, $RW_0\,16$.

Running the rung over the full census: \texttt{CENSUS ... MATCH}, \texttt{STATUS
EXHAUSTED}, $565{,}836{,}800$ placements, 420 survivors, 2{,}520 configurations,
\textbf{zero Costas arrays}, minimum defect 4 attained by 3 configurations,
12.7\,s. Fail-closed merge re-scored all 3 with zero failures:
\texttt{EXHAUSTIVE, NO COSTAS ARRAY FOUND}.

\begin{theorem}[Census-complete rung $k=3$]
\label{thm:rung3census}
No order-32 Costas array contains an order-29 Costas array as a three-row,
three-column deletion minor.
\end{theorem}

This is independent of Theorem~\ref{thm:rungs}' $k=4$ clause rather than stronger
than it (Remark~\ref{rem:independent}): neither implies the other without a
hereditary-minor property that is not established and is not supported by the
order-28 measurements.

\begin{remark}[The two census-complete rungs are independent]
\label{rem:independent}
Deriving the $k=4$ clause from the $k=3$ clause would require every order-29
Costas array to contain an order-28 Costas deletion minor; deriving $k=3$ from
$k=4$ would require the converse extension property. Neither implication is
established. Theorem~\ref{thm:heredity} measures that a single corner deletion
preserves the Costas property only $14.6\%$ of the time at order 28, showing the
required deletion property is far from automatic and therefore providing no basis
for deriving either rung theorem from the other. No counterexample is exhibited;
the point is logical independence of the present proofs, which close different
extension channels --- extension from an order-29 core and extension from an
order-28 core.
\end{remark}

\paragraph{The generator gap, identified.}
The corrected algebraic generator produces 148 of the 164 arrays, and the
shortfall is now explained exactly: the 16 missing arrays form the 2 $RW_0$
equivalence classes produced by the semi-empirical Rickard--Welch method, which
appends a blank row to an exponential $W_1$ array, cyclically shifts, and adjoins
a column. That operation is not corner surgery, and the implemented
corner-deletion/addition closure does not reach its two classes. The gap is
therefore explained by a missing \emph{operation class} rather than by a defect
in the implemented closure; it does not exclude some unrelated genealogy for
those arrays through corner operations from a different seed. The generated set is a verified subset of the
census. This also confirms the order-29 taxonomy: no array of order 29 is
sporadic, but 2 classes lie outside $G\cup PE$ as the generator implements it.

\paragraph{Prior run, superseded.}
The originally recorded run used 120 algebraic cores and reported
$120\times\binom{32}{3}^2\times3!=17{,}713{,}152{,}000$ a priori configurations
in 439\,s, with $2{,}136$ of $2{,}952{,}192{,}000$ core placements surviving the
core-only test; zero Costas arrays, minimum defect 4. The corrected generator's
148 arrays give
$516{,}633{,}600$ placements, 388 survivors, 2{,}328 configurations, zero Costas
arrays, minimum defect 4, 11.6\,s --- the same outcome, now subsumed by
Theorem~\ref{thm:rung3census}.

\begin{theorem}[Rung results]
\label{thm:rungs}
No order-32 Costas array arises as a two-point augmentation of any of the 664
generated order-30 Costas cores. No order-32 Costas array contains an order-29
Costas array as a three-row, three-column deletion minor, and none contains an
order-28 Costas array as a four-row, four-column deletion minor. Among all
$326{,}709{,}248$ two-point augmentations the minimum defect is 2, attained
uniquely up to the dihedral group.

The first conclusion is catalogue-relative: order 30 has never been enumerated,
so its 664 cores are the generated set. The latter two are census-complete,
quantifying over all 164 order-29 and all 712 order-28 Costas arrays.
\end{theorem}

\begin{remark}[Canonical mode is exact]
Restricting cores to one representative per dihedral orbit is exact for
existence and for the minimum defect, because the configuration space is
$D_4$-equivariant and the defect is $D_4$-invariant, so every configuration maps
to one with a canonical core and the same defect. It is \emph{not} exact for
counting solutions.
\end{remark}

\subsection{Genealogy of the score-2 object}
\label{sec:genealogy}

The rung searches identify the parent structure of $N_{32}$ directly. The
order-31 array found at $k=1$ is
\[
C_{31} = (30,2,8,1,14,12,7,4,20,0,18,9,13,21,10,11,26,25,6,28,24,3,15,5,19,22,27,29,16,23,17),
\]
and the order-32 optimum found at $k=2$ is
\[
N_{32}' = (31,30,2,8,1,14,12,7,4,20,0,18,9,13,21,10,11,26,25,6,28,24,3,15,5,19,22,27,29,16,23,17),
\]
which is $C_{31}$ with a single dot adjoined at the corner. $N_{32}'$ is a
dihedral image of the known candidate $N_{32}$.

\begin{observation}[Genealogy]
\label{obs:genealogy}
\[
\text{Welch-30 (algebraic)}
\;\xrightarrow{\ \text{free adjunction}\ }\;
C_{31} \text{ (unique derived order-31 orbit)}
\;\xrightarrow{\ \text{minimum-cost adjunction}\ }\;
N_{32}\ (\text{defect }2).
\]
The first adjunction is free; the second costs 2 and nothing costs less.
\end{observation}

This is a better description of $N_{32}$ than ``Welch-30 plus two dots'': it
identifies an exact intermediate Costas object and localises the failure
entirely to the final adjunction.

\subsection{The toll is a minimum, not a pointwise quantity}
\label{sec:toll}

The bottoms of the two defect spectra coincide after a shift of 2
($8,0,24,16$ appears at defects $0,1,2,3$ for $k=1$ and at $2,3,4,5$ for
$k=2$). This suggested a pointwise parent--child relation. It was tested and is
false.

For each two-dot child $Q$ with defect at most 10, let
$\tau(Q)=\score(Q)-\min\{\score(Q-x),\score(Q-y)\}$ over its two one-dot
parents. Measured distribution:

\begin{center}
\begin{tabular}{@{}lrrrrrrrrrr@{}}
\toprule
$\tau$ & 0 & 1 & 2 & 3 & 4 & 5 & 6 & 7 & 8 & negative \\
\midrule
count & 76 & 48 & 152 & 56 & 296 & 108 & 192 & 24 & 16 & 32 \\
\bottomrule
\end{tabular}
\end{center}

Thirty-two children score \emph{below} both parents, so the defect is not even
monotone along the ladder. The histogram coincidence is a coincidence of
multiplicities.

\begin{observation}
The correct statement is: \emph{the second-rung class has a uniform minimum toll
of 2, attained uniquely up to $D_4$}. It is not the case that each second
adjunction costs exactly 2.
\end{observation}

%=====================================================================
\section{Heredity of the Costas property}
\label{sec:heredity}

The value of deepening the ladder depends on how often a Costas array contains a
smaller one. With a complete census at order 28 this is directly measurable.
Deleting $k$ columns forces deletion of the $k$ rows they occupy; the remaining
dots are compressed and tested.

\begin{center}
\begin{tabular}{@{}rrrrr@{}}
\toprule
$k$ & sub-order & deletions per array & arrays with a Costas sub-array & \\
\midrule
1 & 27 & 28 & 104 / 712 & \textbf{14.6\%} \\
2 & 26 & 378 & 24 / 712 & \textbf{3.4\%} \\
3 & 25 & 3{,}276 & 0 / 712 & \textbf{0.0\%} \\
4 & 24 & 20{,}475 & 0 / 120 sampled & 0.0\% \\
5 & 23 & 98{,}280 & 0 / 120 sampled & 0.0\% \\
\bottomrule
\end{tabular}
\end{center}

\begin{theorem}[Heredity is corner-driven]
\label{thm:heredity}
Over the complete order-28 census, exactly 104 of 712 arrays possess a corner
dot ($14.6\%$, against a random-permutation corner rate of $4/28 = 14.3\%$). Of
the 112 order-27 Costas sub-arrays found, \textbf{104 arise from corner
deletion} and only 8 are non-corner coincidences.
\end{theorem}

The explanation is immediate: corner deletion induces a translation, which
preserves all displacements, whereas any interior deletion induces a compression,
which does not.

\begin{heuristic}[Implication for rung prioritisation]
\label{cor:ladder-value}
If the deletion-minor frequencies at order 32 resemble those measured in the
order-28 census, then an order-32 Costas array would be much more likely to
possess an order-31 or order-30 Costas minor than an order-28 one. This motivates
prioritising small $k$. The $k=1$ and $k=2$ searches are exhaustive relative to
the generated order-30 catalogue; $k=3$ and $k=4$ are census-complete at their
respective core orders. The measurement lowers the expected value of $k\ge5$ but
does not prove those channels empty.
\end{heuristic}

This is the single most consequential measurement in the report. It caps the
reach of branch (ii) of the three-branch decomposition, and correspondingly
enlarges branch (iii).


%=====================================================================
\section{Coordinate classification of the census}
\label{sec:coords}

Costas arrays can be screened by interpolating them as polynomials over a finite
field and reading off the degree. Done naively this measures nothing; done in the
right coordinates it classifies a complete census exactly. Both halves are
recorded here because the failure mode is instructive and was reached twice.

\subsection{Additive-node interpolation has no discriminating power}

The natural first attempt interpolates $h$ against the additive nodes
$0,1,\dots,N-1$ over a prime $p>N$ and screens on the resulting degree, support
size and bidegree rank deficiency. Calibrated on known algebraic arrays:

\begin{center}
\begin{tabular}{@{}lrccc@{}}
\toprule
array & $N$ & $p=67$ & $p=71$ & $p=73$ \\
\midrule
$W_1(13)$ order 12 & 12 & 11/11/0 & 11/11/0 & 11/11/0 \\
$W_1(31)$ order 30 & 30 & 29/29/0 & 29/29/0 & 29/28/0 \\
$W_2(31)$ order 29 & 29 & 28/28/0 & 28/29/0 & 28/29/0 \\
\bottomrule
\end{tabular}
\end{center}

\noindent(degree / support size / maximum rank deficiency over bidegrees
$(d_x,d_y)\le(3,3)$). Maximal degree, full support, zero deficiency --- and
random permutations at $N=29$, $p=67$ over 300 samples give degree 28 in 296
cases and 27 in 4. The known constructions are \emph{more} generic-looking than
random permutations, which occasionally drop a degree by chance.

\begin{observation}[Analytic null]
\label{obs:analyticnull}
For fixed nodes the Newton coefficients of a random permutation vanish
independently with probability about $1/p$, so $P(\deg\le D)\approx
p^{-(N-1-D)}$. Verified: at $N=29$, $p=67$ the observed $P(\deg\le27)=4/300=0.0133$
against the model $1/67=0.0149$. At a threshold $D=18$ this is
$67^{-10}=1.8\times10^{-18}$, and about $5\times10^{-53}$ across three primes and
eight orientations. Monte Carlo estimation of such a null with $10^4$ draws
measures nothing; a reported Clopper--Pearson interval of $[0,\,3\times10^{-4}]$
is some 49 orders of magnitude too wide.
\end{observation}

\subsection{Exponential coordinates}

The Welch relation is $h_i = g^{i+c}-1$, which is linear in the
\emph{multiplicative} parameter $u_i=g^i$, not in the additive index. Screening
against additive nodes destroys exactly the structure sought. Two coordinate
systems suffice for the known families:

\begin{align*}
\text{single-exponential (Welch):}&\quad u_i=g^{i}\in\F_q,\quad v_i=h_i,
  &&h_i=g^{c}u_i-1,\ d=1;\\
\text{dual-exponential (Golomb):}&\quad u_i=\alpha^{i+1},\quad v_i=\beta^{h_i+1},
  &&\alpha^i+\beta^j=1 \Rightarrow v=1-u,\ d=1.
\end{align*}

Offsets are free: $u\mapsto\lambda u$ sends $a_j\mapsto a_j\lambda^{-j}$ and
$v\mapsto\mu v$ scales the interpolant, so neither changes the degree and no grid
search over $c_1,c_2$ is required. Measured:

\begin{center}
\begin{tabular}{@{}lrrl@{}}
\toprule
array & $N$ & min degree & witness \\
\midrule
$W_1(13)$ order 12 & 12 & \textbf{1} & single, $(13,2)$ \\
$W_1(29)$ order 28 & 28 & \textbf{1} & single, $(29,2)$ \\
$W_1(31)$ order 30 & 30 & \textbf{1} & single, $(31,3)$ \\
$G_2(31)$ order 29 & 29 & \textbf{1} & dual, $(31,3,3)$ \\
$G_2(32)$ order 30 & 30 & \textbf{1} & dual, $(32,2,2)$, $\F_{2^5}$ \\
random permutations & 12 & min 8, median 10 & --- \\
random permutations & 29 & min 26, median 27 & --- \\
\bottomrule
\end{tabular}
\end{center}

The analytic null for degree $\le1$ at $N=29$ is $3.5\times10^{-39}$, dominated by
the smallest usable modulus $q=31$ with its 8 primitive roots and 8 orientations:
$64\cdot31^{-27}$.

Extension-field arithmetic is not optional at order 28. The Golomb relation
$q-2=28$ needs $q=30$, not a prime power, so \emph{no intact Golomb array of order
28 exists}; every Golomb-type order-28 array is a corner deletion of $G_2(31)$ or
a double deletion of $G_2(32)$, and the latter lives in $\F_{2^5}$.

\subsection{Complete classification of the order-28 census}

Applying both screens to all 712 arrays, with provenance cross-referenced against
independently generated families:

\begin{center}
\begin{tabular}{@{}lrrlrl@{}}
\toprule
class & arrays & orbits & coordinates & $d$ & provenance \\
\midrule
Welch direct & 672 & 84 & single, $\F_{29}$ & 1 & $W_1(29)$ \\
Golomb, one deletion & 8 & 1 & dual, $\F_{31}$ & 1 & $G_3(31)$ \\
Golomb, two deletions & 24 & 3 & dual, $\F_{2^5}$ & 1 & $G_4(32)$ \\
augmented residue & 8 & 1 & generic in both & $\approx$25 & $G_2(29)$ order 27 $+$ one corner dot \\
\midrule
total & \textbf{712} & \textbf{89} & & & \\
\bottomrule
\end{tabular}
\end{center}

\begin{observation}[Independent agreement with the published classification]
The order-28 enumeration reports its 89 equivalence classes as 84 from $W_1$, 1
from $G_1$, 1 from $G_3$ and 3 from $G_4$. The coordinate screen, which uses no
provenance information whatever, returns 84 single-exponential classes over
$\F_{29}$, 3 dual-exponential over $\F_{2^5}$, 1 dual-exponential over $\F_{31}$
and 1 generic in both. The multisets $\{84,3,1,1\}$ agree exactly, class by
class, from wholly independent derivations.
\end{observation}

\begin{observation}[The emergent $RW_0$ classes are coordinate-generic]
\label{obs:rw0}
Reconstruction of the order-29 census (Section~\ref{sec:catalogue}) removed the
data blockage on this test. Screening the two $RW_0$ equivalence classes --- the
16 arrays produced by the semi-empirical Rickard--Welch method, unreachable by
corner surgery --- gives minimum degree 27 in single-exponential and 26 in
dual-exponential coordinates, for both classes. Genuine arrays of the same order
collapse: $G_2(31)$ to degree 1 dual-exponential, $W_2(31)$ to degree 1
single-exponential. Random permutations of order 29 give minimum 26, median 27.

The emergent arrays therefore sit exactly where random permutations sit and are
\emph{structurally distinct from both algebraic families} under every coordinate
system tested. Their construction is genuine but their coordinate behaviour is
not algebraic in this sense.
\end{observation}

The $W_3(31)$ family is empty at order 28: no double corner deletion of
$W_1(31)$ reaches 28. The generated $G_3(31)$ and $G_4(32)$ families contain
exactly 8 and 24 order-28 arrays respectively, all in the census, matching the
witness split exactly.

\begin{theorem}[Order-28 census classification]
\label{thm:census28}
Every array in the complete order-28 census has degree 1 in single-exponential
coordinates over $\F_{29}$ or in dual-exponential coordinates over $\F_{31}$ or
$\F_{2^5}$, except for a single dihedral orbit of 8 arrays which is generic in
both. That orbit is $G_2(29)$ of order 27 augmented by one corner dot; its
canonical representative is
\[
(0,16,21,12,18,22,23,26,10,7,24,14,27,6,17,5,25,11,20,2,4,19,15,13,8,1,9,3),
\]
and it arises in 8 ways, for instance $G_2(29,\alpha{=}3,\beta{=}3)$ with a corner
dot at $(0,27)$ or $(27,0)$.
\end{theorem}

That orbit is precisely the one the first, deletion-only catalogue generator
missed (Observation~\ref{obs:missing}). The same fact is visible from two
independent directions.

\subsection{Deletion preserves coordinates, addition destroys them}

\begin{observation}
\label{obs:delvsadd}
Corner deletion removes a dot at a boundary of both ranges, so the compression
acting on the survivors is a pure translation, and translations are exactly what
the offset freedom absorbs. Both single and double corner deletion therefore
preserve degree 1. Corner \emph{addition} adjoins a point satisfying no field
relation, and a single off-curve point is sufficient to destroy the
parameterisation.
\end{observation}

The mechanism is exact. If $N-1$ of $N$ points lie on a line $L$ then $P-L$
vanishes at those $N-1$ nodes, so
\[
P(u)-L(u) \;=\; c\prod_{i=1}^{N-1}(u-u_i),
\]
and if the added point lies off $L$ then $c\ne0$ and $\deg P = N-1$.

For the augmented order-28 orbit the parent $G_2(29,3,3)$ has degree 1 in its
exact dual-exponential system $(\alpha,\beta)=(3,3)$. In that same system,
adjoining the corner point raises the child to degree exactly $27=N-1$. The value
26 is attained only after minimising over other $(\alpha,\beta)$ systems at
$q=29$, and is therefore a chance coordinate reduction rather than the degree in
the parent's system; minimising over the full screening family reduces the
reported figure to about 25, exactly as it does for random permutations
($\{24{:}1,\ 25{:}18,\ 26{:}6\}$ over 25 draws). Three distinct quantities are
involved --- 27 in the parent's own system, 26 minimised over $q=29$, about 25
minimised over the full screen --- and they should not share the label ``native
degree''. The residual orbit should be read as ``coordinate-generic''.

This pairs with Theorem~\ref{thm:heredity}: corner deletion preserves the Costas
property $14.6\%$ of the time and, when it fires, preserves the algebraic
coordinates as well; corner addition preserves the Costas property occasionally
but in the augmented order-28 orbit studied here it destroys the parent's
degree-one parameterisation.

\subsection{The intermediate degree band}

A new construction family would appear at degree $2\le d\le N-3$. Screening all
990 dihedral orbits of the complete order-12 census against a matched control of
990 canonical random permutations, field set $\{13,17\}$ identical for both arms:

\begin{center}
\begin{tabular}{@{}lrrrrr@{}}
\toprule
degree & 1 & 7 & 8 & 9 & 10 \\
\midrule
census (990 orbits) & \textbf{12} & \textbf{4} & 45 & 474 & 455 \\
matched random (990) & 0 & 1 & 36 & 475 & 478 \\
\bottomrule
\end{tabular}
\end{center}

Degree $\le7$: census $16/990 = 1.62\%$ against random $1/990 = 0.10\%$. Degree
exactly 1: census $12/990$ against random $0/990$.

\begin{observation}
The degree-1 class is the only statistically resolved separation between the
Costas census and the matched random control.

After removing the degree-1 class, the generic body at $d\in\{8,9,10\}$ has
counts $(45,474,455)$ in the census and $(36,475,478)$ in the control. A Pearson
chi-squared homogeneity test gives $\chi^2=1.454$ on 2 degrees of freedom,
$p=0.483$: the body is statistically indistinguishable from the control under
this test. The extreme lower tail at $d=7$ holds four census orbits against one
control orbit; a two-sided Fisher exact test gives $p=0.215$, so the excess is
not significant. By contrast the degree-1 class holds 12 census orbits against
none, with Fisher $p=4.7\times10^{-4}$.
\end{observation}

\begin{center}
\begin{tabular}{@{}llll@{}}
\toprule
comparison & test & statistic & result \\
\midrule
$d\in\{8,9,10\}$ body & Pearson homogeneity & $\chi^2=1.454$, 2 df & $p=0.483$ \\
$d=7$ lower tail, degree 1 removed & Fisher, two-sided & --- & $p=0.215$ \\
$d=1$ algebraic class & Fisher, two-sided & --- & $p=4.7\times10^{-4}$ \\
\bottomrule
\end{tabular}
\end{center}

\noindent The Fisher comparison for the lower tail treats $d=7$ against
$d\ge8$ within each arm after removing the degree-1 observations, giving the
table $(4,974)$ against $(1,989)$. The comparison for the algebraic class treats
$d=1$ against $d\ge7$ in the two arms, giving $(12,978)$ against $(0,990)$.

The four census orbits attaining $d=7$ do not share a common witness. Two attain
degree 7 only in the single-exponential screen and have no dual-exponential
witness; the other two attain it in distinct dual-exponential systems over
$\F_{13}$, at $(\alpha,\beta)=(11,7)$ and $(\alpha,\beta)=(7,11)$. These are
transposed coordinate choices rather than a parameter system shared by all four,
and no common construction ancestry was detected.

Together with the non-significant lower-tail comparison, this witness diversity
supports reading the four $d=7$ cases as a generic extreme tail rather than a
coherent intermediate-degree family. That is a statement about the population; it
does not classify any individual orbit as random. The population therefore
separates into an exact degree-1 algebraic class and a non-linear remainder whose
bulk is statistically indistinguishable from the matched control.

This closes the inverted search --- enumerating low-degree permutation
polynomials and testing whether they induce Costas arrays --- independently of
Proposition~\ref{prop:polydual}. The window penalty says such a search cannot be
run efficiently at order 32; the measurement above says there is probably nothing
to find.

\begin{remark}[Methodological]
Two distinct analyses in this thread reached conclusions from samples far below
the resolution of the quantity being estimated: $10^4$ Monte Carlo draws against a
null of $5\times10^{-53}$, and roughly 250 candidate polynomials against a window
penalty of $4.15\times10^{-15}$. In both cases the observed count was zero and was
read as evidence of absence. A third analysis drew structural conclusions from
arrays that were not Costas, having bypassed the input validator. The general
lesson is the one already embedded in the rung tooling: assert the invariant
before computing on the object, and compute the null analytically whenever it is
below what sampling can resolve.
\end{remark}

%=====================================================================
\section{Telescoping-sum propagation}
\label{sec:tmprop}

The telescoping identity $\sum_i(h_{i+m}-h_i)=(\text{sum of last }m)-(\text{sum
of first }m)$ couples the endpoint sums of a partial assignment to its
lag-$m$ difference multiset. This section records the inference rule it yields,
a proof that the rule is sound, and its evaluation under a prospectively frozen,
cost-controlled protocol on complete finite populations. It is the one
prospectively gated positive result in this report.

\subsection{Benchmark construction, and a problem with the original gate}
Regenerating the order-12 infeasibility benchmark from scratch --- enumerate all
Costas-valid assignments of the first $s$ columns, then test by exhaustive search
whether each extends to a full order-12 Costas array:

\begin{center}
\begin{tabular}{@{}rrrr@{}}
\toprule
$s$ & Costas-valid prefixes & infeasible & rate \\
\midrule
4 & 10{,}112 & 5{,}120 & 50.63\% \\
5 & 64{,}000 & 56{,}898 & 88.90\% \\
6 & 304{,}288 & \textbf{296{,}656} & 97.49\% \\
7 & 987{,}664 & 979{,}894 & 99.21\% \\
\bottomrule
\end{tabular}
\end{center}

\noindent Among the $304{,}288$ Costas-valid six-leading-column prefixes at
order 12, $296{,}656$ --- about $97.49\%$ --- admit no completion. Local Costas
validity at depth six is therefore a weak predictor of global extendibility. The
benchmark is precisely the complete population of infeasible Costas-valid
six-\emph{leading-column} prefixes at order 12; results on it do not generalise
to arbitrary six-cell partial assignments.

\begin{observation}[The predeclared F3 gate is underpowered]
\label{obs:gate}
The recorded benchmark of ``296 infeasible configurations'' is a \emph{sample},
three orders of magnitude smaller than the $296{,}656$ available at $s=6$. With
$m=296$ draws and an observed union detection rate of $96/296=32.4\%$, the
binomial standard error is
\[
\sqrt{296\cdot0.324\cdot0.676}\;\approx\;8.1,
\]
so the gate at 106 detections sits about $1.25$ standard errors above the
baseline. A propagator with \emph{identical} power would clear it roughly one
time in ten. The gate as frozen cannot distinguish a real improvement from
sampling noise.
\end{observation}

The repair is to run against the complete finite population and preserve the
relative effect size encoded by the original gate. The revised gate therefore
requires
\[
|D_T|\;\ge\;\bigl\lceil 1.104167\,|D_B|\bigr\rceil,
\]
together with zero false positives and a cost ratio of at most two. This is an
exact finite-population comparison, not a significance test: on the complete
population the detection rate is the exact quantity $r_B=|D_B|/296{,}656$, and a
standard error would re-enter only if these prefixes were treated as a sample
from a wider family of orders, depths or geometries. A caveat on scope: the
reconstruction here fixes whole leading columns, and if the original benchmark
fixed arbitrary cells the population will differ. The power calculation does not
depend on which definition is used; a 296-element sample is underpowered either
way.


\subsection{The frozen protocol and the result}
\emph{Audit trail.} An earlier run reported a null failure. It was wrong, through
two implementation defects recorded here as findings rather than footnotes. (i)
The singleton-shaving routine rebuilt each domain from the unary filter on entry,
overwriting the tentative singleton, so shaving never fired; the conclusion that
``shaving contributes nothing'' was an artefact, and the weakened baseline made
the $T_m$ margin vanish. (ii) Relations between two fixed prefix positions were
filtered against the used-difference mask, which contains that pair's own
difference, so every such relation emptied and RA rejected every input. The
second was caught immediately by the complete positive arm. The methodological
conclusion is general: \emph{testing only known-infeasible configurations can
measure detection power but cannot establish even empirical soundness.}

\begin{lemma}[Soundness of the $T_m$ inference rule]
\label{lem:tmsound}
Fix a partial assignment with sound position-value domains $D_i$, and a lag $m$.
Let $k$ be the number of pairs $(i,i+m)$ not wholly inside the prefix, $U$ the
union of the corresponding difference domains, $\det$ the sum of the determined
differences at lag $m$, and $R$ the remaining value set. Put
\[
[\ell,h]=\Bigl[\det+\!\!\sum_{k \text{ smallest of } U}\!\!,\ \
\det+\!\!\sum_{k \text{ largest of } U}\Bigr],
\]
and let $[\tau_-,\tau_+]$ be the corresponding interval for
$T_m=(\text{sum of last }m)-(\text{sum of first }m)$, obtained from the extremal
$f_R$- and $l_R$-subset sums of $R$. If $[\ell,h]\cap[\tau_-,\tau_+]=\varnothing$
then the partial assignment has no Costas completion.
\end{lemma}
\begin{proof}
Let $\sigma$ be any Costas completion. Telescoping gives the identity
$\sum_{i=0}^{n-1-m}\bigl(\sigma_{i+m}-\sigma_i\bigr)
=\sum_{j=m}^{n-1}\sigma_j-\sum_{i=0}^{n-1-m}\sigma_i
=\bigl(\textstyle\sum\text{last }m\bigr)-\bigl(\sum\text{first }m\bigr)=T_m(\sigma)$,
valid for every sequence. Since $\sigma$ is Costas, its $n-m$ differences at lag
$m$ are pairwise distinct, so the $k$ undetermined ones are $k$ distinct integers,
each lying in its domain and hence in $U$ by soundness of the domains. Any $k$
distinct elements of a finite set have sum between that of its $k$ smallest and
its $k$ largest, so the left side lies in $[\ell,h]$. The unknown parts of the two
endpoint sums are subset sums of $R$ of fixed sizes, each between the extremal
subset sums of that size, so $T_m(\sigma)\in[\tau_-,\tau_+]$. As the two sides are
equal, the intervals meet. Contrapositive. Both interval constructions are
outer approximations: they may include sums incompatible with the individual
variable domains or with joint row usage, but they exclude no sum realised by a
completion, so disjointness is a valid certificate and can never be manufactured.
\end{proof}

\subsubsection*{Benchmark}\noindent $\mathcal B_{n,s}$, every Costas-valid
$s$-leading-column prefix at order $n$, labelled by exhaustive completion. At
$(12,6)$: $304{,}288$ valid, $296{,}656$ infeasible, $7{,}632$ completable,
independently reproduced by a second completion procedure.
\textbf{Baselines.} $D_{\mathrm{RA}}$, relation-algebra closure by path
consistency on difference-relation masks under sumset composition; $D_{L_3}$,
position--value domains with unary filtering, arc consistency, Hall matching and
singleton shaving; $D_B=D_{\mathrm{RA}}\cup D_{L_3}$.
\textbf{Candidate.} $D_T=D_B$ plus the rule of Lemma~\ref{lem:tmsound}, so
$D_B\subseteq D_T$ and the gate reduces to marginal coverage.
\textbf{Gate, frozen before the run:} $|D_T|\ge\lceil1.104167|D_B|\rceil$, zero
false positives, cost ratio $\le2$ on the \emph{declared deterministic propagation-work
counter}: \texttt{workRA} increments per relation initialisation, per sumset
shift and per composition; \texttt{workL3} per unary-filter test, per
arc-consistency support test and per matching call; \texttt{workTM} per domain
scan entry and per interval test. It is a frozen instrumentation metric, not a
literal instruction count or a hardware-independent measure of total cost.

\begin{theorem}[F3, computational]
\label{thm:f3}
Among the $296{,}656$ infeasible Costas-valid six-leading-column prefixes of order
12, the frozen baseline $D_B$ detects $77{,}074$. Adding the $T_m$-coupled rule
detects $117{,}880$, including $40{,}806$ additional infeasible prefixes, at
$1.050$ times the baseline declared propagation-work count. Neither system
rejects any of the $7{,}632$ completable prefixes.
\end{theorem}

The frozen threshold was $\lceil1.104167\times77{,}074\rceil=85{,}103$; the
candidate exceeds it by $32{,}777$ detections, an improvement of
$40{,}806/77{,}074\approx52.94\%$ against a required $10.4167\%$, for about $5\%$
additional work. \textbf{The gate passes on all three criteria.}

\medskip\noindent External validation on further complete populations (the formal
gate was declared at $(12,6)$; these rows are exploratory generalisation):

\begin{center}
\begin{tabular}{@{}rrrrrrl@{}}
\toprule
$n$ & $s$ & $|D_B|$ & required & $|D_T|$ & cost & verdict \\
\midrule
10 & 5 & 3{,}882 & 4{,}287 & 7{,}388 & 1.060 & pass \\
10 & 6 & 47{,}230 & 52{,}150 & 54{,}104 & 1.017 & pass \\
11 & 6 & 88{,}652 & 97{,}887 & 119{,}734 & 1.023 & pass \\
11 & 7 & 325{,}922 & 359{,}873 & 341{,}520 & 1.008 & gate inapplicable$^\ast$ \\
12 & 7 & 875{,}236 & 966{,}407 & 969{,}828 & 1.006 & pass \\
\bottomrule
\end{tabular}
\end{center}

\noindent$^\ast$\emph{Ceiling saturation}, not failure: with
$|D_B|=325{,}922$ and $|\mathcal B|=342{,}010$ the maximum possible relative
improvement is $4.94\%$, so a $10.4167\%$ gate is arithmetically impossible. A
multiplicative total-coverage gate is inapplicable whenever the baseline rate
exceeds $1/1.104167\approx90.6\%$. For such regimes the appropriate secondary
criterion is miss reduction,
$\rho_{\mathrm{miss}}=|D_T\setminus D_B|/|\mathcal B\setminus D_B|$; here
$15{,}598/16{,}088\approx96.95\%$ of the baseline's remaining misses are
eliminated. This metric is recommended for future high-baseline benchmarks and is
\emph{not} used retroactively --- $(12,6)$ passed its own frozen gate.

\begin{observation}[Mechanistic reading]
\label{obs:f3}
The $T_m$ rule supplies substantial marginal propagation beyond the implemented
RA and $L_3$ procedures across multiple complete leading-prefix benchmarks: at the
declared benchmark it improves detection by about $52.94\%$ over baseline for
about $5.0\%$ additional deterministic work. The $40{,}806$ marginal detections
are certificates the implemented relation-algebra closure and $L_3$ singleton
shaving do not produce, so the contribution cannot be explained as a different
execution order for those two procedures. On the tested leading-prefix
benchmarks, RA contributes no detections beyond $L_3$
($D_{\mathrm{RA}}\setminus D_{L_3}=\varnothing$ at every point), so the working
stack there is $L_3+T_m$; this is benchmark-specific and does not generalise.
\end{observation}

\noindent\textbf{What F3 establishes:} a correctly implemented new propagation
mechanism, sound by Lemma~\ref{lem:tmsound} and exhaustively validated on the
complete positive arm; substantial marginal detection beyond the frozen baseline;
reproducibility across several complete finite populations; a favourable
coverage-to-cost ratio. \textbf{What it does not:} improved total runtime in a
full order-32 search; the same marginal gain at order 32; that RA is always
subsumed by $L_3$; or that no other domain-consistency method derives the same
contradictions --- that would need a formal expressiveness theorem.



\subsection{F3$'$: does the rule reduce an exact search?}
\label{sec:f3prime}

F3 measured detection on static prefixes. Whether those deductions reduce the
tree of a running solver is a different question, and it was asked under a
separately frozen protocol. One executable selects the propagator by flag; every
non-propagation choice --- branching order, value order, symmetry handling,
incremental displacement check, thread count --- is identical across arms.
Relation algebra is omitted from the candidate, since it contributed no
detections beyond $L_3$ on every tested benchmark. The invocation schedule
(propagate when at most 8 positions are free) was \emph{declared a priori and
never tuned}, so no order is a development order and all three are held out.

\begin{center}
\begin{tabular}{@{}rrrrrrrrr@{}}
\toprule
$n$ & solutions & $N_I$ & $N_{L_3}$ & $N_{L_3+T_m}$ & $R_{\text{tree}}$ & $W$ ratio & $t_{L_3}$ & $t_{L_3+T_m}$ \\
\midrule
11 & 4{,}368 & 1{,}437{,}514 & 379{,}290 & 259{,}414 & $0.316$ & $1.055$ & 5.0\,s & 5.9\,s \\
12 & 7{,}852 & 6{,}977{,}361 & 1{,}522{,}289 & 1{,}119{,}809 & $0.264$ & $1.047$ & 31.8\,s & 36.8\,s \\
13 & 12{,}828 & 33{,}831{,}238 & 5{,}785{,}132 & 4{,}175{,}962 & $0.278$ & $1.037$ & 174.4\,s & 197.1\,s \\
\bottomrule
\end{tabular}
\end{center}

\noindent Solution counts match the published censuses. Set identity is
established \emph{exactly}, not by hashing: each mode emits every solution in
canonical form, and the sorted files are compared byte for byte. At orders 11,
12 and 13 all three modes produce identical files
(\texttt{sha256} prefixes \texttt{3b346690b48c2527}, \texttt{2baed7d03cc6f9de},
\texttt{b6cf5cfc1580bff2} for 4{,}368, 7{,}852 and 12{,}828 solutions). The two
order-independent 64-bit hashes are retained as a fast regression guard;
aggregate hashes are not collision-free certificates of set equality, so the
claim rests on the exact comparison. The fixed column-by-column DFS also
generates each full assignment at most once, so duplicates are impossible by
construction.

\begin{observation}[F3$'$ verdict: structural pass, operational failure]
\label{obs:f3prime}
Against the gate frozen before the run: \textbf{(a)} identical solution sets ---
pass; \textbf{(b)} $N_{L_3+T_m}\le0.90N_{L_3}$ --- pass, measured $0.684$,
$0.736$, $0.722$; \textbf{(c)} $W_{L_3+T_m}\le1.25W_{L_3}$ --- pass, measured
$1.055$, $1.047$, $1.037$. The \emph{scientific} gate passes at every order.
\textbf{(d)} $t_{L_3+T_m}<t_{L_3}$ on a majority --- \textbf{fails}, 0 of 3; the
candidate is $12$--$18\%$ slower in wall time. The \emph{operational} gate fails.
$T_m$ removes a quarter to a third of the remaining search tree for about $4\%$
additional deterministic work, and the current implementation does not convert
that into elapsed time.
\end{observation}

The larger finding is about the baseline, not the candidate. At order 13,
$L_3+T_m$ explores only $12.3\%$ of the incremental-only tree
($4{,}175{,}962$ against $33{,}831{,}238$ nodes) and yet takes $197$\,s against
$3.1$\,s. Across the three orders incremental-only explores $5.5$ to $8.1$ times
more nodes and runs $53$--$63$ times faster, because the exact bitmask displacement
check costs almost nothing per node while shaving costs a great deal. At these
orders heavyweight propagation is an operational loss \emph{in either
configuration}, and $T_m$ inherits that. The node-reduction trend is stable
rather than growing with order ($0.316$, $0.264$, $0.278$), so nothing here
supports projecting a crossover; establishing one would require either a far
cheaper $L_3$, a sparser invocation schedule, or orders well beyond those
enumerable here.

The honest reading: $T_m$ is a correct and genuinely useful deduction rule that
pays for itself in \emph{deductions per unit of declared work} and does not pay
for itself in seconds against this baseline.


\subsection{The representation boundary}
\label{sec:repboundary}

Lemma~\ref{lem:tmsound} is proved for any value of the formal contract. It does
not follow that \texttt{f3\_complete.cpp} builds one. That gap is now an object
in the development rather than a paragraph in this report.

\paragraph{Proved.} The bitmask/list bridge. The executable stores position
domains as \texttt{unsigned D[u]}, bit $k$ meaning ``the $k$-th remaining value is
still available at position $u$''; the Lean development uses lists of indices.
\texttt{domainList} and \texttt{mem\_domainList} identify the two exactly, and
\texttt{mem\_domainList\_of\_submask} shows that every filtering step --- each is
\texttt{D[u] \&= \~{}(1u <\!< k)}, producing a submask --- only shrinks a domain.
Pool membership proved against the \emph{filtered} domains therefore reads back
as membership in the original.

\paragraph{Not proved.} \texttt{CppObligations} is a structure whose eight fields
are the complete residue, each documented against the construct it refers to:

\begin{center}
\begin{tabular}{@{}ll@{}}
\toprule
obligation & claim about the C++ \\
\midrule
\texttt{domains\_sound} & \texttt{l3Detect} clears a bit only for a value no completion can take \\
\texttt{pairs\_enumerated} & the lag loop visits each unresolved pair once; \texttt{a < s} classifies it \\
\texttt{pools\_match} & the two \texttt{push\_back} branches generate the two pool functions \\
\texttt{dedup\_matches} & \texttt{sort}/\texttt{unique} has the membership of \texttt{dedupInt} \\
\texttt{prefix\_sums\_extremal} & \texttt{pre\_lo}, \texttt{pre\_hi} are the extremal $k$-subset sums \\
\texttt{endpoints\_match} & \texttt{slo}, \texttt{shi}, \texttt{tlo}, \texttt{thi} are the interval endpoints \\
\texttt{cardinalities\_correct} & \texttt{fk}, \texttt{fR}, \texttt{lk}, \texttt{lR} are the true window sizes \\
\texttt{no\_overflow} & no accumulation overflows; all sums are $O(n^2)$ \\
\bottomrule
\end{tabular}
\end{center}

\noindent An opaque \texttt{Prop} field, however, is an unfalsifiable
placeholder: it can be neither proved nor refuted, and a differential harness can
only ever refute. \texttt{Reference.lean} therefore supplies a Lean reference
implementation of the computation \texttt{tmDetect} performs and discharges the
shallow obligations as theorems about it:

\begin{center}
\begin{tabular}{@{}ll@{}}
\toprule
former obligation & now \\
\midrule
\texttt{pools\_match} & \texttt{ref\_fixedFree\_eq}, \texttt{ref\_freeFree\_eq} (by \texttt{rfl}, no axioms) \\
\texttt{dedup\_matches} & \texttt{ref\_dedup\_preserves}, \texttt{ref\_dedup\_nodup} \\
\texttt{prefix\_sums\_extremal} & \texttt{ref\_lower\_eq} (by \texttt{rfl}, no axioms) \\
\texttt{endpoints\_match} & \texttt{ref\_disjoint\_iff} \\
\texttt{no\_overflow} & \texttt{ref\_sum\_bounded}: sum bounded by $\text{length}\times n$, so $O(n^2)$ \\
\texttt{pairs\_enumerated}, \texttt{cardinalities\_correct} & absorbed below \\
\bottomrule
\end{tabular}
\end{center}

\noindent What survives is \texttt{CppResidue}, with \textbf{two} fields rather
than eight. \texttt{domains\_sound} --- that \texttt{l3Detect} clears a bit only
for a value no completion can take --- is the deep one, and is not a
representation question at all: it asserts that arc consistency and singleton
shaving are correct filters. \texttt{cpp\_matches\_reference} is a single
refinement claim about straight-line code, and is precisely what a differential
harness \emph{can} test on the complete order-12 population.

\begin{remark}
The reduction matters more than the count. Eight opaque propositions were an
unbounded liability; two typed contracts, one of them testable and the other a
named mathematical property of a filtering algorithm, are a work item. Replacing
placeholders with a reference implementation is what made the difference ---
five obligations became theorems because there was finally something to prove
them about.
\end{remark}


\subsection{Order reachability settles the taxonomy without the names}
\label{sec:reachability}

The work package above seeks to resolve the ambiguous $G_3$/$G_4$/Taylor/Lempel
designations into exact source-defined equations before specialising each to
$\F_{32}$. That resolution is worth having, but the question it exists to answer
can be settled first, and independently of every naming decision, by asking which
\emph{orders} corner surgery reaches rather than what each family is called.

Every $W_1(p)$ for prime $p\le43$ and every $G_2(q)$ for supported prime power
$q\le43$ was generated over all primitive parameter pairs, and each seed set was
closed under corner deletion and corner addition:

\begin{center}
\begin{tabular}{@{}lrrl@{}}
\toprule
family & native order & seeds & orders reachable by corner surgery \\
\midrule
$G_2(29)$ & 27 & 144 & 26, 27, 28, 29 \\
$G_2(31)$ & 29 & 64 & 28, 29, 30 \\
$G_2(32)$ & 30 & 180 & 28, 29, 30 \\
$G_2(37)$ & 35 & 144 & \textbf{34, 35} \\
$W_1(29)$ & 28 & 12 & 26, 27, 28, 29 \\
$W_1(31)$ & 30 & 8 & 29, 30 \\
$W_1(37)$ & 36 & 12 & \textbf{36} \\
\bottomrule
\end{tabular}
\end{center}

\begin{observation}[The corner-surgery order gap]
\label{obs:ordergap}
Corner surgery from these seeds reaches exactly
$\{26,\dots,30\}\cup\{34,35,36\}$. \textbf{Orders 31, 32 and 33 are reached by
nothing.} In particular:
\begin{itemize}
\item $G_2(37)$ reaches 34 by one corner deletion and stops. No second deletion
  survives, so whatever equations the $G_4$ and $G_5$ labels denote, the families
  they name are \emph{empty at $q=37$}. Since 37 is the only prime power that
  could reach 33 or 32 by deletion, the conditional constructions cannot produce
  either order.
\item $W_1(37)$ admits no corner deletion at all, so order 35 is not reached that
  way either.
\end{itemize}
\end{observation}

This makes the naming question moot for the order-32 problem. Whatever exact
equations $G_3$, $G_4$, $G_5$, Taylor and Lempel denote, each is a corner
deletion of a $G_2$ or $W_1$ seed, and no such deletion has order 31, 32 or 33.
The taxonomy remains worth completing for correctness of the record and for
neighbouring-order work, but it cannot supply an order-32 witness.

\begin{remark}[Order 31 is not a corner-surgery descendant]
The gap is informative about order 31 as well. Order-31 Costas arrays exist, and
rung $k=1$ exhibits eight of them in one dihedral orbit --- but they do not appear
in the table above, because rung $k=1$ performs \emph{general} one-point
augmentation into a compressed sub-board, not corner addition. So the eight
order-31 arrays are precisely non-corner augmentations, which is consistent with
Theorem~\ref{thm:heredity}: corner operations preserve structure and are rare,
and order 31 needs the operation that does not. Applying the same general
augmentation twice is rung $k=2$, exhausted over all 664 generated order-30 cores
with zero order-32 arrays (Theorem~\ref{thm:rungs}).
\end{remark}

%=====================================================================
\section{Structural and algebraic approaches assessed}
\label{sec:structural}

This section records approaches that were proposed, implemented and closed. Each
subsection states the proposal, the closure, and the quantitative evidence.
Observation~\ref{obs:uniformity} is the common cause in most cases; where a
different mechanism is responsible it is identified.

\subsection{Modular decomposition modulo 16}
\label{sec:mod16}

\paragraph{Proposal.} Since the heights are a permutation of 32 consecutive
integers, every residue modulo 16 occurs exactly twice. A complete modulo-16
template is a partition of the 32 columns into 16 labelled pairs, and each pair
admits one orientation bit, so lifting a completed template costs only
$2^{16}=65{,}536$ trials. Search balanced templates, prune with the modular
difference capacities $n_{d,0}\le 2$, $n_{d,r}\le 4$, then lift.

\paragraph{Closure.} The decomposition is exact but is a change of variables, not
a reduction. The template space has size $32!/2^{16}\approx 4\times10^{30}$ and
the lift costs $2^{16}$; the product is exactly $32!$.

More decisively, the modular constraint is a relaxation of what the exact search
already enforces. Exact difference equality implies modular equality, so every
constraint the modular layer can fire is a weakening of one the exact solver
fires, and it fires later, since most residue pairs become comparable only once
both columns are placed. The capacity bounds are not extra inference either: if
five distance-$d$ edges share a residue modulo 16, two must share an exact
difference, which the exact search detects as soon as both endpoints are
assigned. \emph{The capacity constraints are a patch for information the
abstraction discarded, and a patch cannot recover more than the discard cost.}

\paragraph{Evidence.} The original validation at $(N,M)=(8,4)$ recovered all 444
order-8 Costas arrays from 2{,}520 balanced templates, of which 1{,}936 survived
the capacity filter --- a reduction of only $\mathbf{23.2\%}$. See also
Section~\ref{sec:2adic}, where the modulo-16 capacity is shown to recover 8\% of
the true collision count at order 32.

The associated lift experiment is nevertheless a clean negative result worth
retaining: all $2^{16}$ orientations of the fixed modulo-16 template of $N_{32}$
were enumerated, the minimum defect is 2, it is attained once, and the
second-best defect is 13. The isolation has an algebraic explanation rather than
a modular one: flipping an orientation bit shifts a height by $\pm16$, half the
range, destroying the circular structure on which the Welch core depends.

\subsection{Relation algebra}
\label{sec:ra}

\paragraph{The precise algebra.} Atoms are exact signed differences; the working
algebra is $\Cm(\Z)$ (or $\Cm(\Z^2)$ carrying the redundant horizontal
coordinate). Converse is negation. Composition is the group addition table, so
composition of domains is the \emph{sumset}. Row indices are not part of the atom
type: relations are homogeneous, and the board data lives at the network level.

Properties, all automatic and independent of Costas: $\Cm(\Z^2)$ is complete,
atomic, \emph{integral} ($1'$ is an atom), \emph{commutative}, not symmetric,
and every atom is functional ($a;a^\smile = a^\smile;a = 1'$). It is
\emph{representable} by the regular representation, hence lies in
$\RRA = \bigcap_{n\ge4}\RA_n$ unconditionally. Its consistency relation is
\emph{rigid}: $C(a,b,c)$ holds iff $c=a+b$, so each pair of atoms has exactly one
consistent completion --- the extreme opposite of the flexible-atom property that
makes Ramsey relation algebras work.

\begin{proposition}[Truncation destroys associativity]
Restricting composition to a box $D_m=[-(m-1),m-1]^2$ does not yield a relation
algebra. With $X=\{(31,0)\}$, $Y=\{(1,0)\}$, $Z=\{(-31,0)\}$ one has
$(X;Y);Z=\emptyset$ because $(32,0)\notin D_{32}$, while $X;(Y;Z)=\{(1,0)\}$.
Bounded intervals are not closed under addition, so this fails for every $m$.
\end{proposition}

\begin{proposition}[A Costas array is not a representation]
Restricted to the box there are $63^2-1 = 3968$ non-identity atoms but only
$32\times31 = 992$ ordered pairs on 32 nodes. Every representation --- square,
relativised or weak --- is injective with $h(0)=\emptyset$, so every atom
requires a nonempty image. No representation of any kind exists on 32 points,
for reasons independent of the Costas condition.
\end{proposition}

What a Costas array \emph{is}, is a board-relative injective network. Because
each atom is functional, the triangle law holds with equality, so the network is
determined by the star at any one node --- this is precisely the vertex-potential
reduction, and it is the entire dividend the rigidity pays.

\begin{proposition}[Corrected board-network characterisation]
\label{prop:board}
The axioms identity, converse, cocycle, $d=j-i$, $|e|\le m-1$, strictness and
off-diagonal injectivity do \emph{not} characterise Costas arrays: strictness is
vacuous once $d=j-i$ is imposed. The missing axiom is that every off-diagonal
label has \emph{nonzero vertical component}. Counterexamples satisfying all the
stated axioms with a repeated height, normalised to $h_0=0$:
\begin{center}
\begin{tabular}{@{}lrrrrrrr@{}}
\toprule
$m$ & 2 & 3 & 4 & 5 & 6 & 7 & 8 \\
\midrule
count & 1 & 12 & 110 & 852 & 7{,}520 & 61{,}688 & 528{,}506 \\
\bottomrule
\end{tabular}
\end{center}
\end{proposition}

\paragraph{Expressivity.} Relation-algebraic reasoning is essentially the
three-variable fragment of first-order logic. The Ramsey correspondence works
because ``no monochromatic triangle'' is a three-point condition. The Costas
condition is a four-point condition on $(i,i{+}d,j,j{+}d)$ and lies outside that
fragment; the cocycle move evades this by promoting edge colours to variables,
but the resulting algebra is a group complex algebra and is representable, so no
obstruction survives. The natural home for a four-point condition is Maddux's
$\RA_n$ hierarchy or cylindric algebras of dimension $\ge4$.

\paragraph{Measured strength of RA closure.} The closure operation
$D_{ik}\leftarrow D_{ik}\cap(D_{ij}+D_{jk})$ to fixpoint, together with Régin
GAC on the AllDifferent constraint of each horizontal-distance class, was
implemented on 63-bit edge masks and benchmarked against vertex-domain
propagation on 300 shared Costas-valid six-cell configurations at $n=12$, of
which 296 are globally infeasible:

\begin{center}
\begin{tabular}{@{}lrr@{}}
\toprule
propagator & refuted & unsound \\
\midrule
RA closure (edge domains + distance AllDifferent) & 43 & 0 \\
vertex-domain L2 (height GAC + difference-class GAC + channelling) & 25 & 0 \\
vertex-domain L3 (L2 + pairwise extension) & 83 & 0 \\
\midrule
$\text{RA}\setminus\text{L2}=21$, $\text{L2}\setminus\text{RA}=3$ & &\\
$\text{RA}\setminus\text{L3}=13$, $\text{L3}\setminus\text{RA}=53$ & &\\
union(RA, L3) & 96 & \\
\bottomrule
\end{tabular}
\end{center}

\begin{observation}
RA closure is \emph{not} the vertex-domain CP model renamed. It is edge-domain
propagation, incomparable with vertex-domain propagation, contributing 13 of 296
refutations the latter misses and missing 53 it catches. An explicit instance
refuted by RA closure and not by L3 is $n=12$, cells
$\{(3,2),(4,5),(5,9),(6,6),(8,11),(11,7)\}$; the largest single composition step
is $D_{0,1}\leftarrow D_{0,1}\cap(D_{0,8}+D_{8,1})$, removing the atoms
$\{-11,+11\}$. RA refutes nothing at configuration size $\le5$ (0 of 600 sampled
at each of $s=4,5$).
\end{observation}

\paragraph{Detection by configuration size ($n=12$, complete catalogue).}
\begin{center}
\begin{tabular}{@{}lrrrrrr@{}}
\toprule
$s$ (fixed cells) & 1 & 2 & 3 & 4 & 5 & 6 \\
\midrule
RA closure & 0 & 0 & 0 & 0 & 0 & $10.26\%$ \\
vertex L2 & 0 & 0 & 0 & 0 & 0 & $7.18\%$ \\
vertex L3 & 0 & 0 & 0 & 0 & 0 & $27.18\%$ \\
\bottomrule
\end{tabular}
\end{center}
At $n=13$, L2 detects $1.36\%$ at $s=6$: detection \emph{falls} as $n$ grows.

\paragraph{Pointed instances.} A single fixed cell yields nothing. Exhaustive
sweeps over all cells at orders 25, 26, 27 and 29 (2{,}871 pointed instances)
refuted \emph{none}, and the vertex-domain propagator pruned exactly
$\text{cells}\times(n-1)$ values, i.e.\ nothing beyond removing the fixed value
from the other domains. RA closure at order 27 pruned 70{,}200 edge atoms over
729 cells, an average of 96 out of 36{,}504 per instance ($0.26\%$), and refuted
none. The order-32 central-row representatives behave identically.

\paragraph{Score-2 prefix diagnostic.} Supplying prefixes of $N_{32}$ of
increasing length, closure refutes nothing at lengths 2 through 6; at length 7
the collision $(d,\Delta)=(1,1)$ between columns $0\to1$ and $5\to6$ becomes
explicit and is caught by the disequality itself. \emph{Closure gives no early
warning whatsoever}: detection coincides exactly with explicit violation.

\paragraph{Structural reason.} For distance class $d$ there are $n-d$ difference
variables drawing from $2n-2$ values, so occupancy is at worst $1/2$. The
AllDifferent constraints are never within a factor of two of tight, so Hall and
matching filtering have nothing to bite on until domains have already collapsed.
This is not an implementation weakness. Improving propagation independently
within each sparse distance-class AllDifferent is unlikely to change this
behaviour; additional power must come from coupling information \emph{across}
classes, as the $T_m$ result of Section~\ref{sec:tmprop} demonstrates.

\subsection{Matroid and circuit structure}
\label{sec:matroid}

\paragraph{Proposal.} Dias da Silva and Nathanson prove that if $X$ is a finite
$B_{2h-1,h-1}$-subset of an abelian group, the $B_h$-subsets of $X$ form the
independent sets of a matroid, and in particular all maximal ones have the same
cardinality. For DDCs, $h=2$, so the hypothesis is that the grid is a
$B_{3,1}$-set.

\paragraph{Closure.} The grid is not a $B_{3,1}$-set at any order: already in the
$3\times3$ grid, $\{(0,0),(1,1),(1,2)\}$ and $\{(0,1),(0,2),(2,0)\}$ are disjoint
with common sum $(2,3)$. And the conclusion fails, not merely the hypothesis.
Inclusion-maximal DDCs in the $n\times n$ grid have different cardinalities:

\begin{center}
\begin{tabular}{@{}ll@{}}
\toprule
$n$ & sizes of inclusion-maximal DDCs \\
\midrule
4 & 4, 5, 6 \\
5 & 5, 6, 7, 8 \\
6 & 6, 7, 8, 9 \\
7 & 6, 7, 8, 9, 10 \\
\bottomrule
\end{tabular}
\end{center}

For permutation DDCs it is worse. In the $8\times8$ grid,
$\{(1,4),(2,6),(3,3),(5,1),(6,5)\}$ is a partial Costas configuration with
\emph{zero} one-dot extensions, while order-8 Costas arrays have 8 dots. The
minimum inclusion-maximal partial Costas size runs $3,4,4,5,5$ for $n=4,\dots,8$
against maxima $4,\dots,8$.

There is therefore no exchange property, hence no rank function, no circuit
axioms, no greedy and no basis exchange.

\paragraph{Why.} The affine matroid of $\Z^2$ has rank 3 after homogenisation, so
circuits have support at most 4. The Costas condition forbids exactly those with
coefficient vector $(1,-1,-1,1)$ (parallelogram) and $(1,-2,1)$ (three
equidistant collinear points). It does \emph{not} forbid $(3,-1,-1,-1)$,
$(1,2,-1,-2)$ or any other affine dependency. Costas selects a unimodular
sub-family of circuits by coefficient pattern, and toric machinery --- Graver
bases, circuit ideals, Gr\"obner degree bounds, normality --- is closed under the
full circuit set and cannot see such a restriction.

\subsection{Affine and toric invariants}
\label{sec:affine}

\paragraph{Complete linear invariant space.} Let $\chi(h)_{d,e}=\#\{i:
h_{i+d}-h_i=e\}$ be the occupancy vector. Exhausting $S_n$ gives the complete
space of linear functionals constant on all permutations --- that is, every
moment-type obstruction that can exist at the linear level:

\begin{center}
\begin{tabular}{@{}rrrr@{}}
\toprule
$n$ & ambient dimension & affine hull & invariant dimension \\
\midrule
4 & 18 & 11 & 7 \\
5 & 32 & 21 & 11 \\
6 & 50 & 37 & 13 \\
7 & 72 & 55 & 17 \\
\bottomrule
\end{tabular}
\end{center}

Three families were identified: (I) row sums $\sum_e\chi_{d,e}=n-d$; (II) signed
column sums $\sum_d(\chi_{d,e}+\chi_{d,-e})=n-|e|$; and (III) the telescoping
symmetry $T_d = T_{n-d}$, where $T_d=\sum_e e\,\chi_{d,e}$, proved by
$T_d=\Sigma - S_d - S_{n-d}$ with $S_k$ the prefix sum. A residual family of
size $\lfloor (n-1)/2\rfloor$ was not identified in closed form. Family (II) is
the exact affine content: the vertical difference multiset over all ordered pairs
is fixed by $n$ alone, so only its distribution across horizontal classes is
free.

\paragraph{Feasibility.} Setting $y_{d,g}=\chi_{d,g}+\chi_{d,-g}\in\{0,1,2\}$,
families (I) and (II) become a $31\times31$ matrix with row sums and column sums
both $(31,30,\dots,1)$, each totalling $\sum_{k=1}^{31}k=496$. Such a matrix was
constructed greedily, so the affine relaxation is satisfiable at $n=32$ and gives
no obstruction. The totals are equal \emph{identically in $n$}, which is the
signature of a constraint system balanced by construction.

\paragraph{Toric reduction.} Reducing displacements modulo $n$, the $n(n-1)$
ordered differences land in $(\Z_n\setminus\{0\})^2$, which has $(n-1)^2$ cells.
Since $n(n-1)>(n-1)^2$ there are at least $n-1$ forced coincidences.

\begin{corollary}
No permutation DDC is a Sidon set in $\Z_n\times\Z_n$; equivalently, no Costas
array is doubly periodic. At $n=32$: 992 ordered differences into 961 cells,
forced excess exactly 31.
\end{corollary}

This is a genuine toric obstruction, but to double periodicity, not existence,
and it is uniform in $n$. The one place toric structure has produced a real
conclusion in the literature is the single periodicity of $W_1$, which is why
there can be no forbidden positions at orders $n=p-1$.

\subsection{The 2-adic tower}
\label{sec:2adic}

\paragraph{Proposal.} Replace the flat modulo-16 model by a filtration
$\Z/2\leftarrow\Z/4\leftarrow\Z/8\leftarrow\Z/16\leftarrow\Z/32$, with
capacities $2^{5-s}$ per residue class at level $s$, propagating nested modular
domains $R^{(s)}_{ij}\subseteq \Z/2^s$.

\paragraph{The capacity claim is false.} It is off by a factor of 2, because it
treats modulo 32 as the exact level. Exact differences lie in $[-31,31]$, which
is 63 values, so the injective level is modulo \textbf{64}. The correct signed
capacities $N_q(r)$ are:

\begin{center}
\begin{tabular}{@{}lrr@{}}
\toprule
modulus & residue 0 & other residues \\
\midrule
2 & 30 & 32 \\
4 & 14 & 16 \\
8 & 6 & 8 \\
16 & \textbf{2} & \textbf{4} \\
32 & 0 & 2 \\
64 & 0 & 1 \\
\bottomrule
\end{tabular}
\end{center}

The claimed ``modulo-16 capacity two'' rule is refuted by $N_{32}$ itself: at
$d=1$, residue 12 modulo 16 carries three edges with exact differences
$-20,-4,12$, all distinct and all liftable; and the original analysis records a
saturated class $(d,r,n)=(5,1,4)$ with four edges. Adopting the rule would reject
$N_{32}$ and every genuine Costas array with a saturated class.

\begin{theorem}[Domination]
\label{thm:domination}
Let $\pi_s:[-31,31]\to\Z/2^s$ be reduction, $D_{ij}$ the exact edge domains and
$R^{(s)}_{ij}$ the tower domains. Since $\pi_s(A+B)\subseteq\pi_s(A)+\pi_s(B)$
--- strictly whenever an integer sum leaves the board or two lifts merge --- and
the level-$s$ capacity is the image of the exact AllDifferent, every value the
tower removes is also removed by flat exact propagation. Hence tower refutes
$\Rightarrow$ flat refutes, never the converse.
\end{theorem}

\paragraph{Cost.} The premise that coarse levels are cheap is false at $n=32$:
the exact domain is 63 values and already fits in one 64-bit word, whereas the
full tower needs $2+4+8+16+32+64=126$ bits and six compositions instead of one.
The tower is strictly more expensive and strictly weaker. Its leaf count
telescopes exactly:
\[
\binom{32}{16}\binom{16}{8}^2\binom{8}{4}^4\binom{4}{2}^8\binom{2}{1}^{16}
= 32!,
\]
so, like the flat modulo-16 model, it is a change of variables.

\paragraph{Measured pruning power.} Define the sharpest defect bound the
level-$s$ capacities can give, $L_s=\sum_{d,r}\max(0, m_{s,d,r}-N_s(r))$.

\begin{center}
\begin{tabular}{@{}rrrr@{}}
\toprule
$s$ & modulus & $C_s$ for $N_{32}$ & $L_s$ for $N_{32}$ \\
\midrule
1 & 2 & 2490 & 0 \\
2 & 4 & 1185 & 0 \\
3 & 8 & 544 & 0 \\
4 & 16 & 210 & 0 \\
5 & 32 & 56 & 1 \\
6 & 64 & 2 & \textbf{2} \\
\bottomrule
\end{tabular}
\end{center}

$L_6$ equals the exact defect identically, because modulo 64 with capacity 1
\emph{is} the Costas condition. Everything below is a strict relaxation, and
$L_4=0$: the modulo-16 level cannot distinguish a defect-2 near-miss from a
genuine Costas array.

Over 20{,}000 random order-32 permutations (mean defect 90.39):

\begin{center}
\begin{tabular}{@{}rrrrr@{}}
\toprule
$s$ & modulus & fraction with $L_s>0$ & mean $L_s$ & $\mathrm{corr}(L_s,\score)$ \\
\midrule
2 & 4 & $0.84\%$ & 0.01 & 0.04 \\
3 & 8 & $44.3\%$ & 0.82 & 0.20 \\
4 & 16 & $99.7\%$ & 7.45 & 0.40 \\
5 & 32 & $100\%$ & 28.19 & 0.67 \\
6 & 64 & $100\%$ & 90.39 & 1.00 \\
\bottomrule
\end{tabular}
\end{center}

The $99.7\%$ figure is misleading: random permutations have mean defect 90, so
anything catches them. The operative number is $L_4=0$ for the one permutation
of interest.

\subsection{XOR and Boolean Hensel lifting}
\label{sec:xor}

\paragraph{Proposal.} At level $s$, a modular collision
$h_{i+d}-h_i\equiv h_{j+d}-h_j \pmod{2^s}$ with integer discrepancy
$\Delta^{(s)}$ separates at level $s+1$ exactly when
\[
\epsilon_{i+d}\oplus\epsilon_i\oplus\epsilon_{j+d}\oplus\epsilon_j
= 1 \oplus (\Delta^{(s)}\bmod 2),
\]
giving a pure 4-XOR system amenable to Gaussian elimination in CryptoMiniSat.

\paragraph{The algebra is correct; imposing it is unsound.} Separating at level
$s+1$ is sufficient for exact separation but never necessary: a pair may
legitimately stay congruent for several levels, requiring only
$\vtwo(D)<\infty$ for $D$ the difference of differences, not $\vtwo(D)=s$.
Demanding the XOR at each level demands $\vtwo(D)$ equal the level at which the
collision first appears.

\begin{center}
\begin{tabular}{@{}lrrrrrr@{}}
\toprule
$\vtwo(D)$ & 0 & 1 & 2 & 3 & 4 & 5 \\
\midrule
order-31 parent $C_{31}$ & $51.1\%$ & $25.7\%$ & $12.6\%$ & $6.7\%$ & $3.1\%$ & $0.9\%$ \\
Welch-30 core & $49.6\%$ & $26.3\%$ & $13.3\%$ & $6.9\%$ & $3.3\%$ & $0.7\%$ \\
\bottomrule
\end{tabular}
\end{center}

About $49\%$ of same-distance pair-pairs in a genuine Costas array do not
separate at level 1. A hard XOR system at that level rejects $C_{31}$ and the
Welch core --- arrays known to exist. For a non-existence programme this is the
fatal failure mode: spurious UNSAT, invisible to any amount of UNSAT running.

The correct constraint is a disjunction (separate now, or defer and carry the
pair forward), which makes the XOR right-hand side a free variable; the system is
then not pure XOR and the Gaussian-elimination argument evaporates.

\paragraph{What survives.} Read as a \emph{filter} rather than a constraint, the
identity is sound and useful: if the four bit-0 values XOR to 1, that pair is
separated forever. Roughly half die at each level, so after five levels only the
$\vtwo(D)\ge5$ pairs remain --- 39 of 4{,}495 on $C_{31}$. Whether this beats a
single 64-bit hash insertion per pair is a benchmark question.

\subsection{Finite-field coordinate-order families}
\label{sec:gf32}

\paragraph{Proposal.} Search $h(i)=\rho_{\mathrm{out}}\!\left(a\,
[\rho_{\mathrm{in}}^{-1}(i)]^k+b\right)$ over $\F_{32}$, with $\rho$ ranging over
structured coordinate orderings.

\paragraph{Expanded exhaustive search.} With 240 orderings (120 polynomial-basis
bit permutations $\times$ \{binary, reflected Gray\} traversal) chosen
independently on input and output, the family has
$30\times31\times32\times240^2 = 1{,}714{,}176{,}000$ members. Quotienting $b$ by
the vertical reflection $h\mapsto 31-h$ (realised as $b\mapsto b\oplus
c_{\mathrm{out}}$ with $c_{\mathrm{out}}=31$ for binary output and 16 for Gray,
since $\mathrm{gray}(31-i)=\mathrm{gray}(i)\oplus16$) leaves $857{,}088{,}000$ to
sweep. The earlier best of 50 is reproduced exactly, confirming the
parameterisation is a faithful superset.

\begin{center}
\begin{tabular}{@{}lrrr@{}}
\toprule
range & tested & survived $\le55$ & best \\
\midrule
$k\in[1,16)$ & 428{,}544{,}000 & 40{,}762 & 41 \\
$k\in[16,31)$ & 428{,}544{,}000 & 49{,}057 & \textbf{37} \\
\bottomrule
\end{tabular}
\end{center}

Best: defect 37 at $k=26$, $a=17$, $b=12$, in$(\sigma=119,\text{binary})$,
out$(\sigma=72,\text{Gray})$.

\paragraph{Matched comparison.} 857{,}088{,}000 uniform random permutations at the
identical cutoff give best 41. Tail counts at exactly matched sample size:

\begin{center}
\begin{tabular}{@{}lrrrrrrrrrrrr@{}}
\toprule
defect & 37 & 39 & 40 & 41 & 43 & 45 & 47 & 49 & 51 & 53 & 55 & $\le55$ \\
\midrule
family & 1 & 1 & 3 & 4 & 25 & 84 & 367 & 1267 & 4417 & 13237 & 37372 & 89{,}819 \\
random & 0 & 0 & 0 & 1 & 2 & 6 & 42 & 211 & 613 & 2143 & 6702 & 14{,}887 \\
\bottomrule
\end{tabular}
\end{center}

Body quantiles over 3{,}000{,}000 matched draws each:

\begin{center}
\begin{tabular}{@{}lrrrrrrr@{}}
\toprule
& mean & p0.01 & p0.1 & p1 & p5 & p25 & p50 \\
\midrule
family & 102.70 & 55 & 60 & 66 & 72 & 80 & 87 \\
random & 90.33 & 59 & 64 & 70 & 75 & 84 & 90 \\
\bottomrule
\end{tabular}
\end{center}

$D_4$ deduplication on 3{,}000{,}000 draws gives 2{,}985{,}196 distinct orbits,
so the tail is not a duplication artefact. Adding input translations (family
size $5.5\times10^{10}$, 400{,}000{,}000 random draws) gives best 40 and 85 hits
at $\le45$ against 81 predicted by scaling the base family: no gain.

\begin{observation}[Constant factor, not slope]
\label{obs:slope}
Both tails decay at approximately $1.8\times$ per defect point. The family's
advantage is a constant multiplicative factor of about 6, worth roughly three
defect points, and it does not compound. Extrapolating at the measured slope,
reaching defect $\le25$ needs $1.8^{12}\approx1.2\times10^3$ times the family
(about $10^{12}$ members) and reaching 0 needs $1.8^{37}\approx2.8\times10^9$
times the family (about $2.4\times10^{18}$). The entire expanded family is
$1.7\times10^9$ and is now exhausted.
\end{observation}

\begin{remark}[Methodological lesson]
The gate that was applied disjunctively --- ``reach defect $\le25$ \emph{or}
show a clear lower-tail shift'' --- is flawed, because any structured family
beats random by a constant factor. A successor gate should specify a change in
the tail \emph{slope}. Better still, the diagnostics that separate a promising
family from a merely enriched one are: multiplicity and orbit count of the
minimum, the gap above the minimum, persistence under family enlargement, and
the existence of a structural parent--child explanation. The augmentation family
of Section~\ref{sec:ladder} passes all of these (minimum 2, multiplicity 8, one
orbit, defect 3 empty, explicit parent); the $\F_{32}$ family passes none.
\end{remark}


%=====================================================================
\section{The dyadic defect filtration}
\label{sec:dyadic}

The mod-64 injectivity noted in Section~\ref{sec:2adic} can be turned into a
filtration on the defect set rather than on the search. This section records what
that yields: one unconditional theorem, and a measurement showing why the
associated optimisation problem does not bear on the existence question.

\subsection{Exact counting and the filtration}

Directed triples $(r,\ell,h)$ correspond to $0\le a<b<c\le31$ via $a=r$,
$b=r+\ell$, $c=r+\ell+h$, so there are $\binom{32}{3}=4960$ directed defect
conditions. The diagonal count is $\sum_{\ell=1}^{15}(32-2\ell)=240$, and after
the duality quotient $R_{\ell,h,r}=R_{h,\ell,r}$ there are
$240+\tfrac12(4960-240)=2600$ distinct defect expressions.

Since $-62\le R_{\ell,h,r}\le62$ at order 32, $R=0$ in $\Z$ if and only if
$64\mid R$. Writing $Z_k(p;\Omega)=\#\{\omega\in\Omega : 2^k\mid R_\omega(p)\}$,
one has $Z_1\ge Z_2\ge\cdots\ge Z_6$ and
\[
p \text{ is Costas} \iff Z_6(p;\Omega_{32})=0 .
\]
With $E_k=\#\{\omega : 2^k\,\|\,R_\omega\}$ the escape recurrence is
$Z_{k+1}=Z_k-E_k$.

\subsection{An unconditional lower bound}

\begin{theorem}[Diagonal parity floor]
\label{thm:parityfloor}
For the diagonal defects $D_{\ell,r}=p_{r+2\ell}-2p_{r+\ell}+p_r$ at order 32,
let $e_0=\#\{i \text{ even}: p_i \text{ even}\}$. Then
\[
Z_1^{\mathrm{diag}} \;=\; 2\binom{e_0}{2}+2\binom{16-e_0}{2}
\;=\; 2(e_0-8)^2+112 \;\ge\; 112,
\]
with equality exactly when $e_0=8$.
\end{theorem}
\begin{proof}
$2\mid D_{\ell,r}$ iff $p_r\equiv p_{r+2\ell}\pmod2$, and $r$ and $r+2\ell$ have
the same index parity, so the 240 diagonal pairs are exactly the same-parity
index pairs. Among the 16 even indices, $e_0$ carry even values; since there are
16 even values in total, the odd indices carry $16-e_0$ even values and $e_0$ odd
ones. Counting same-value-parity pairs within each index class gives the stated
expression.
\end{proof}

This is the only unconditional universal bound the investigation has produced.
Verified identically against 400 random permutations, with the floor 112
attained.

\subsection{Why the mod-4 optimisation does not bear on existence}

Let $q_i=p_i\bmod4$, let $Z_2^{\mathrm{diag}}(q)$ count diagonal triples with
$4\mid q_r+q_{r+2\ell}-2q_{r+\ell}$, and let $m_2$ be its minimum over balanced
residue words, of which there are $32!/(8!)^4 = 99{,}561{,}092{,}450{,}391{,}000$.
A stochastic search produced a balanced word with $Z_2^{\mathrm{diag}}=16$,
independently re-evaluated over all 240 triples, so $0\le m_2\le16$; the exact
value is open.

The statistic is invariant under the lift, since it depends only on
$q=p\bmod 4$. Calibrating it against arrays that actually exist:

\begin{center}
\begin{tabular}{@{}lrr@{}}
\toprule
population & $Z_2^{\mathrm{diag}}$ & fraction of triples \\
\midrule
order-32 minimisation witness & 16 of 240 & $6.7\%$ \\
order-28 Costas census (712 arrays) & 32--56, mean 48.6 of 182 & $26.7\%$ \\
order-28 random permutations (3000) & 28--71, mean 44.2 of 182 & $24.3\%$ \\
order-32 random permutations (2000) & min 39, mean 58.4 of 240 & $24.3\%$ \\
\bottomrule
\end{tabular}
\end{center}

\begin{observation}
\label{obs:m2}
$Z_2^{\mathrm{diag}}$ does not separate Costas arrays from random permutations:
the complete order-28 census sits at mean 48.6 against a random mean of 44.2 on
the same 182 triples, slightly \emph{above} random and well inside the spread.
Real Costas arrays do not avoid level-2 survivors; they let them escape at higher
levels, as the full profiles show --- Welch-30 runs
$Z_1\ldots Z_6=[102,50,25,9,2,0]$ and the order-31 parent $C_{31}$ runs
$[111,57,20,7,1,0]$. Minimising $Z_2^{\mathrm{diag}}$ therefore moves away from
the region where Costas arrays live, and no order-28 Costas array arises from a
word with $Z_2^{\mathrm{diag}}$ anywhere near the minimum.
\end{observation}

Determining $m_2$ exactly is a well-posed finite problem. It does not bear on
whether an order-32 Costas array exists, because the quantity being minimised
carries no Costas signal.

\subsection{The clique reduction and $m_2\ge1$}

The scalar minimisation can in fact be settled, and the structure does it
without a solver. Every same-parity index pair $(i,k)$ with $k-i=2\ell$ satisfies
$1\le\ell\le15$ automatically, so \emph{the 240 diagonal triples are exactly the
240 same-parity index pairs}, each with a determined midpoint $j=(i+k)/2$.
Survival requires $u_i=u_k$ and $v_i\oplus v_k=u_i\oplus u_j$, so for fixed $u$
the zero-survivor condition forces
\[
v_i\oplus v_k \;=\; 1\oplus u_i\oplus u_j
\]
on \emph{every} pair with $u_i=u_k$. The constraint graph is therefore four
\emph{complete} graphs, indexed by (index parity) $\times$ ($u$-value), with
midpoint-generated edge labels $\lambda_{ik}=1\oplus u_i\oplus u_{(i+k)/2}$.

\begin{remark}[A tempting argument that does not work]
It is tempting to conclude inconsistency from equation counting: a clique on $s$
vertices carries $\binom{s}{2}$ equations in $s$ unknowns. That is invalid. The
$\F_2$ incidence matrix of $K_s$ has rank $s-1$, and the remaining
$\binom{s-1}{2}$ equations are cycle conditions which may all hold. Taking
$\lambda_{ik}=w_i\oplus w_k$ for arbitrary $w$ makes every equation consistent at
any size: verified for $s=8,12,16$ (28, 66 and 120 equations respectively, all
consistent). The inconsistency below is a property of the \emph{midpoint}
labelling, not of the redundancy.
\end{remark}

\begin{proposition}[Computational]
\label{prop:m2}
For every $u\in\F_2^{32}$ and every $v\in\F_2^{32}$, the mod-4 word $q=u+2v$ has
at least one diagonal defect divisible by 4. Hence $m_2\ge1$, and combined with
the verified score-16 witness, $1\le m_2\le16$.
\end{proposition}

An exhaustive parity union-find over $u$ (with $u_0=0$ by translation) visits 744
nodes, reaches index 17, and finds no consistent $v$. Two remarks on scope. The
search imposes \emph{no} balance condition on $v$ and, as a separate run
confirms, imposing balance on $u$ changes nothing: with and without
$\sum_i u_i=16$ the search is byte-identical at 744 nodes and zero solutions. The
statement is therefore stronger than $m_2\ge1$ --- no $v$ exists at all, balanced
or not, for any $u$.

This is a \emph{computational} proposition, not a structural theorem. The
obstruction first becomes unavoidable on positions $0\ldots17$: for every
smaller prefix some $u$ is consistent, and at $m=18$ all $2^{17}=131{,}072$ words
fail. There is accordingly no short case analysis, and a structural proof --- that
the midpoint labelling always yields a frustrated triangle --- remains open and
would be worth more than the exact value of $m_2$.

The exact object is
\[
m_2 \;=\; \min_{\substack{u\in\F_2^{32}\\ \sum_i u_i=16}}
\operatorname{fr}_{\mathrm{bal}}(G_u),
\]
the balance-constrained frustration index of the four midpoint-labelled complete
signed graphs, where the minimisation over $v$ must respect
$\sum_{i:u_i=0}v_i=\sum_{i:u_i=1}v_i=8$. Frustration index on $\le16$ vertices is
$2^{16}$ colourings computed exactly, so the residual work is the outer search
over $u$.

\subsection{The survivor hypergraph and the escape system}
\label{sec:escape}

Scalar counts discard which defects share which positions. Retaining it, choose
bit $b_k$ and note $R\mapsto R+2^kA_{k,\omega}$ with
$A_{k,\omega}=b_{r+\ell+h}-b_{r+\ell}-b_{r+h}+b_r$. A survivor escapes iff
$A_{k,\omega}\equiv R_\omega/2^k+1 \pmod 2$, so killing every survivor in one
step is an affine system over $\F_2$,
\[
H_k b_k = c_k,
\]
with rows indexed by $S_k$, columns by positions, $H_k[\omega,i]$ the parity of
the multiplicity of $i$ in the rectangle, and $c_k[\omega]=R_\omega/2^k+1$. Let
$\Lambda_k(p)$ count solutions that also respect the residue balance the lift
requires. Calibrated at order 28, 60 Costas arrays against 60 random
permutations over all 3{,}276 rectangle defects:

\begin{center}
\begin{tabular}{@{}rlrrrr@{}}
\toprule
$k$ & arm & mean $|S_k|$ & mean rank & inconsistent & $\Lambda_k>0$ \\
\midrule
1 & Costas & 1623.3 & 25.0 & 60/60 & 0 \\
1 & random & 1644.3 & 25.0 & 60/60 & 0 \\
2 & Costas & 756.6 & 25.0 & 60/60 & 0 \\
2 & random & 825.6 & 25.0 & 60/60 & 0 \\
3 & Costas & 351.3 & 25.0 & 60/60 & 0 \\
3 & random & 416.4 & 25.0 & 60/60 & 0 \\
\bottomrule
\end{tabular}
\end{center}

\begin{observation}
\label{obs:escape}
The one-step all-escape system is inconsistent for \emph{every} permutation
tested in both arms, and $\Lambda_k\equiv0$. The statistic separates nothing
because it is vacuous: with over 1{,}600 survivors at level 1, a rank-25 system
in 28 unknowns is over-determined by two orders of magnitude, and genuine Costas
arrays do not satisfy it either. They do not kill their survivors in one step;
they need only $S_6=\varnothing$ and reach it progressively. The one visible
difference, that Costas arrays carry slightly fewer survivors at levels 2 and 3,
is the expected consequence of $S_6=\varnothing$ rather than an independent
signal.
\end{observation}

This is the stopping criterion the dyadic programme set for itself: if Costas
arrays and random permutations are indistinguishable under rank, consistency and
balanced-lift statistics, the branch closes. They are, and it does. \subsection{The corrected statistic does separate}
\label{sec:muk}

Demanding $s=\mathbf0$ was the error; the exact transition is
$H_k b = s+\mathbf1+\kappa_k$ with a free survivor mask $s$, subject to
$\sum_{i:a_i=t}b_i = 2^{4-k}$. The right statistic is then not whether all
survivors can be killed at once, but how aggressively the actual lift reduces
them. Fixing $p \bmod 2^k$ and re-lifting with 250 random balanced bit vectors,
at order 28 over 20 Costas arrays and 20 random permutations:

\begin{center}
\begin{tabular}{@{}rlrrrrr@{}}
\toprule
$k$ & arm & $|S_k|$ & actual $|S_{k+1}|$ & sampled min & sampled mean & actual percentile \\
\midrule
1 & Costas & 1628.2 & 759.8 & 766.9 & 817.8 & \textbf{4.0} \\
1 & random & 1653.2 & 825.1 & 780.5 & 830.9 & 45.9 \\
2 & Costas & 764.5 & 350.1 & 347.2 & 386.0 & \textbf{4.7} \\
2 & random & 837.4 & 426.1 & 381.1 & 421.8 & 56.4 \\
\bottomrule
\end{tabular}
\end{center}

\begin{observation}
\label{obs:muk}
Costas arrays lift near-optimally. Their actual bit sits at roughly the fourth
percentile of balanced re-lifts of the same base, and at $k=1$ their mean
next-level count (759.8) is below the mean \emph{minimum} over 250 random
balanced re-lifts (766.9). Random permutations sit at the median, 45.9 and 56.4.
This is a genuine topology-sensitive separation of the kind the all-escape test
failed to produce.
\end{observation}

Two cautions before this is treated as leverage. First, it may be tautological:
a Costas array must reach $S_6=\varnothing$, so its filtration is obliged to
decay faster than a generic lift's, and the percentile may be measuring only that
obligation. The test that distinguishes the cases is whether \emph{near}-Costas
permutations, of defect 1 or 2, also lift at the low percentile; if they do, the
statistic is predictive and could drive a search heuristic, and if only exact
Costas arrays do, it is a restatement of the terminal condition. Second, the
signal is a percentile, not a bound, so it yields a heuristic rather than a sound
prune: an array may lift badly at one level and recover.

What survives unconditionally is Theorem~\ref{thm:parityfloor} and
Proposition~\ref{prop:m2}, both exact and both about the filtration rather than
about existence.

%=====================================================================
\section{Partition counting catalogue}
\label{sec:partition}

This section consolidates the counting arguments available at the root, with the
verdict for each.

\subsection{Local palette per separation}

At separation $m$ the class needs $L_m=32-m$ distinct signed values from a
palette of $62$, leaving $30+m$ unused. Occupancy is
$(32-m)/62$, maximal at $m=1$ where it is exactly $50\%$; every other class is
sparser. \emph{This is why independent AllDifferent and Hall filtering per class
achieves nothing.}

\subsection{Global magnitude budget}

For $g\in\{1,\dots,31\}$ exactly $32-g$ unordered value pairs have absolute
difference $g$, giving
\[
\sum_{m=1}^{31}\bigl(\chi_{m,g}+\chi_{m,-g}\bigr) = 32-g .
\]
The resulting $31\times31$ matrix $y_{m,g}\in\{0,1,2\}$ has both margins equal to
$(31,30,\dots,1)$, each totalling 496. Evaluating the Hall condition over every
prefix set, the tightest is the whole set with demand 496 and availability 496:
\emph{slack zero}, but that zero is the accounting identity, not a constraint.

\begin{remark}[Correction]
It is nevertheless \emph{not} true that the magnitude budget is permanently
vacuous. After domain contraction an edge may realise only a restricted set of
magnitudes, and a proper subset can then develop a Hall deficiency even though
the totals agree globally. The correct statement is: the magnitude margins do
not bind at the root, but may bind after adaptive domain contraction.
\end{remark}

\subsection{The telescoping sum rule}
\label{sec:sumrule}

This is the only root-level constraint omitted from earlier catalogues that does
real work.

\begin{proposition}
Let $k=\min(m,32-m)$. Then
\[
T_m \;=\; \sum_i \Delta_{m,i}
      \;=\; \sum_{\text{last }k} h_i - \sum_{\text{first }k} h_i,
\qquad |T_m| \le m(32-m),
\qquad T_m = T_{32-m}.
\]
Moreover $T_1 = h_{31}-h_0 \ne 0$, and $T_{16} = 496 - 2\sum_{i<16}h_i$ is
\emph{even}.
\end{proposition}

Counting rainbow sets that are also sum-admissible, by exact dynamic programming
over all $\binom{62}{L_m}$ subsets:

\begin{center}
\begin{tabular}{@{}rrrrrr@{}}
\toprule
$m$ & $L_m$ & rainbow sets & sum-admissible & surviving & bits \\
\midrule
1 & 31 & $4.654\times10^{17}$ & $1.524\times10^{17}$ & $32.74\%$ & 1.611 \\
2 & 30 & $4.509\times10^{17}$ & $2.669\times10^{17}$ & $59.2\%$ & 0.76 \\
3 & 29 & $4.099\times10^{17}$ & $3.157\times10^{17}$ & $77.0\%$ & 0.38 \\
4 & 28 & $3.496\times10^{17}$ & $3.074\times10^{17}$ & $87.9\%$ & 0.19 \\
6 & 26 & $2.098\times10^{17}$ & $2.038\times10^{17}$ & $97.1\%$ & 0.04 \\
16 & 16 & $2.733\times10^{14}$ & $1.367\times10^{14}$ & $50.0\%$ & 1.000 \\
$\ge$12 (others) & & & & $\approx100\%$ & $\approx0$ \\
\bottomrule
\end{tabular}
\end{center}

The $m=1$ row excludes sum zero ($+0.024$ bits over the naive figure); the
$m=16$ row is the parity constraint, worth exactly one bit.

\begin{observation}
The sum of marginal log-reduction factors across all 31 classes is
$\mathbf{4.10}$ bits, against $\log_2 32! = 117.7$. This is a sum of
\emph{marginal} reductions, not a joint information measure: rows $m$ and $32-m$
share one $T$, so at most about 16 of the 31 filters are independent. All the
reduction sits at $m=1,2,3$ and $16$.
\end{observation}

\subsection{Residue, parity and valuation partitions}

The signed residue capacities $N_q(r)$ are tabulated in
Section~\ref{sec:2adic}. Among the 62 signed values, 32 are odd and 30 even, so
only the $m=1$ class ($L_1=31>30$) is forced to contain at least one odd
difference. Modulo 4 the capacities are $14,16,16,16$. The 2-adic valuation
partition is exact and nested:
\[
\#\{\vtwo=0\}=32,\;\;
\#\{\vtwo=1\}=16,\;\;
\#\{\vtwo=2\}=8,\;\;
\#\{\vtwo=3\}=4,\;\;
\#\{\vtwo=4\}=2,\qquad \text{total }62.
\]
No valuation class is forced at any separation.

\subsection{Empirical occupancy}

Over the complete order-28 census (712 arrays, palette 54):

\begin{center}
\begin{tabular}{@{}rrrrr@{}}
\toprule
$m$ & $L_m$ & values ever realised & unseen & $L_m/\text{seen}$ \\
\midrule
1--18 & 27 \dots 10 & \textbf{54} & 0 & 0.50 $\to$ 0.19 \\
19 & 9 & 52 & 2 & 0.17 \\
22 & 6 & 50 & 4 & 0.12 \\
24 & 4 & 50 & 4 & 0.08 \\
27 & 1 & 48 & 6 & 0.02 \\
\bottomrule
\end{tabular}
\end{center}

For every separation up to 18 the \emph{entire} palette is realised somewhere in
the census. Over the 729 magnitude cells, $60.5\%$ sit at $y=0$, $27.1\%$ at
$y=1$ and $12.4\%$ at $y=2$; only 11 cells ($1.5\%$) are never occupied, all at
$m\ge19$ where the class has fewer than ten edges. There is no static forbidden
colour phenomenon.

\subsection{Verdict}

\begin{center}
\begin{tabular}{@{}lp{0.55\textwidth}@{}}
\toprule
constraint & status \\
\midrule
per-class palette & never binds; maximum occupancy $50\%$ \\
magnitude budget, Hall & balanced by construction; root slack 0 is an identity;
  may bind after contraction \\
\textbf{sum rule $T_m$} & \textbf{binds at $m=1,2,3,16$; 4.10 marginal bits of 117.7} \\
residue capacities & necessary, dominated by exact propagation; $8\%$ recovery at
  modulo 16 \\
parity, modulo 4 & one forced odd at $m=1$; nothing else \\
2-adic valuation & exact partition; no forced class \\
\textbf{diagonal parity floor} & \textbf{the only entry that binds
  unconditionally: $Z_1^{\mathrm{diag}}\ge112$ (Thm.~\ref{thm:parityfloor})} \\
empirical palette & fully realised to $m=18$ \\
\bottomrule
\end{tabular}
\end{center}

\emph{No uniform partition bites, because every class is at most half full.}

%=====================================================================
\section{An anti-Ramsey displacement bound}
\label{sec:antiramsey}

The rigidity of the Costas atom structure (Section~\ref{sec:ra}) blocks the
Ramsey correspondence, but the dual --- rainbow, anti-Ramsey --- framing yields a
usable bound for the augmentation problem.

\begin{theorem}[$\ell_\infty$ collision bound]
\label{thm:linfty}
Place a Costas core $K$ of order $n$ on a $B\times B$ board by deleting $k=B-n$
columns and $k$ rows. Two core pairs with compressed displacement vectors
$(\delta,\epsilon)$ and $(\delta',\epsilon')$ can collide on the board only if
\[
\sgn(\epsilon)=\sgn(\epsilon'), \qquad
|\delta-\delta'| \le k, \qquad
|\epsilon-\epsilon'| \le k .
\]
\end{theorem}

\begin{proof}
The column expansion is strictly increasing with $c_i-i\in[0,k]$ non-decreasing,
so the board column distance satisfies $d\in[\delta,\delta+k]$. The row expansion
gives $e\in[\epsilon,\epsilon+k]$ when $\epsilon>0$ and $e\in[\epsilon-k,\epsilon]$
when $\epsilon<0$: the sign is preserved and the magnitude only grows. Two
intervals of length $k$ meet only under the stated conditions, and intervals of
opposite sign never meet.
\end{proof}

\paragraph{Measured.} For an order-28 core at $k=4$, candidate colliding
pair-pairs drop from $\binom{378}{2}=71{,}253$ to $5{,}678$, a $12.5\times$
reduction. Verified empirically over $798{,}453$ observed collisions across
20{,}000 random placements: \emph{zero} violations.

\paragraph{Consequence for search.} The bound does not speed up the order-28 rung,
because the per-placement test buckets 378 displacements by column distance and
aborts on the first collision after roughly ten hash insertions, which beats
iterating 5{,}678 candidates. Its payoff is the \emph{row-side DFS} that the same
monotonicity licenses: choosing deleted rows in increasing order makes a growing
prefix of core pairs testable. Pairs decidable after fixing only the first
deleted row at position $p_1$:

\begin{center}
\begin{tabular}{@{}lrrrrrrr@{}}
\toprule
$p_1$ & 4 & 8 & 12 & 16 & 20 & 24 & 28 \\
\midrule
pairs decidable (of 378) & 6 & 28 & 66 & 120 & 190 & 276 & 378 \\
\bottomrule
\end{tabular}
\end{center}

Measured over six sampled column deletions at $k=4$: $215{,}760$ flat nodes
versus $34{,}842$ DFS nodes, a $\mathbf{6.2\times}$ reduction. This would matter
for $k=5,6,7$ --- but see Heuristic~\ref{cor:ladder-value}, which removes the
motivation for those rungs.

%=====================================================================
\section{Satisfiability and distributed solving}
\label{sec:sat}

\subsection{Why SAT is the wrong tool for enumeration here}

Three reasons, which compound.

The dedicated search performs the exact incremental Costas check for every
newly instantiated displacement at very low cost: the bitmask DFS checks
every newly created displacement with one 64-bit bit test, which \emph{is} exact
propagation, at $4.3$--$14.5$ million nodes/s. A CDCL solver reaches the same
pruning through unit propagation over an $O(n^4)$-clause encoding with watched
literals, clause database management and restarts.

Conflict clauses do not generalise: the constraint system is a stack of mutually
unconstraining AllDifferents over difference classes, the regime in which
resolution has exponential lower bounds.

The encoding loses the structure that makes the problem tractable: the
vertex-potential formulation collapses 496 edge variables to 31 heights, and CNF
over dot variables discards that.

\begin{observation}[Calibration]
A distributed SAT system whose Costas model-counting case study reaches order 17
is matched by a single core in about 14 minutes (Section~\ref{sec:baseline}). The
gap between a SAT case study at order 17 and the dedicated enumeration at order
29 is twelve orders at the measured $4.80\times$ per order, about
$2\times10^{8}$ in search effort. That is not recovered by adding cores.
\end{observation}

\subsection{Where distributed SAT does earn its place}

Three places: independent certification of a sample of branches from a dedicated
search; bounded-neighbourhood questions where certificates matter most, such as
extending a radius-$t$ exclusion around $N_{32}$; and cross-implementation
replication of rung branches. The scale ceiling is real: cube-and-conquer has
certified the Pythagorean triples problem (about 4 CPU-years, 200 TB of proof)
and Schur number five (about 14 CPU-years, 2 PB), the two largest mathematical
proofs produced. Exhaustive order 32 at $10^4$--$10^5$ CPU-years is
$10^3$--$10^4$ times larger, with an exabyte-scale certificate. Cube-and-conquer
is the right architecture at the wrong scale.

\subsection{The generated instances}

Two families were produced with a verified encoder. Variables $x_{i,j}$ for a dot
at column $i$ row $j$; exactly one per row and column; auxiliaries $y_{d,e,i}$
with $x_{i,j}\wedge x_{i+d,j+e}\to y_{d,e,i}$ (only that direction is needed);
and at most one $i$ per class $(d,e)$. Defect $\le1$ relaxes this to at most two
per class, a per-class collision flag, and at most one flag true, via Sinz
sequential counters. At $n=32$: about 32k variables and 864k clauses for the
defect-0 form, 97k and 1.02M for defect $\le1$.

\paragraph{Encoder verification.} Every permutation of $\{0,\dots,N-1\}$ was
checked against the CNF at $N=6,7,8$; CNF satisfaction coincides exactly with
being a Costas array in all $720$, $5{,}040$ and $40{,}320$ cases. The
biconditional (census) encoding forces all 240 auxiliaries at $N=7$ with zero
free. The radius (Sinz) encoding agrees with brute force on all 5{,}040
permutations at $N=7$, $r=2$. A census run at $N=12$ produced 7{,}852 models, all
Costas, no duplicates, matching the published census.

\paragraph{The lemma instance.} Selector variables choose where the new column and
row insert into $C_{31}$, with 1024 ternary clauses pinning the resulting array.
UNSAT at defect $\le1$ certifies that every one-point augmentation of $C_{31}$
has defect at least 2. This instance is decided independently by 1024
brute-force evaluations (minimum defect 2, attained twice), so it functions as a
\emph{calibration instance} rather than a mathematical result.

\subsection{The DAG decomposition, and what it does}

A node is a subset of the CNF's clauses named by index. An edge $a\to b:V$ means
solve $a$, and carry each model's values on $V$ to constrain $b$. A chain of
compatible solutions through to the terminal node solves the whole formula, which
is why the terminal node owns every clause. Parallelism comes from sibling nodes
and, dominantly, from \emph{fan-out}: if node $a$ has $M$ models, node $b$ is
instantiated $M$ times as independent subproblems. Effectiveness turns entirely
on the size of the edge variable set (the separator) and the number of upstream
models; the standard decomposition tool uses min-fill, the classical treewidth
heuristic.

Both decompositions here are hand-built, because the Costas primal graph is dense
--- every pair of columns is coupled through some $y$ variable --- so there is no
small separator for min-fill to find. The lemma uses a two-node generator/worker
split with a 64-variable separator and fan-out 1024. The neighbourhood uses a
cumulative prefix chain over column blocks, where the separator is the
\emph{entire prefix}: displacement constraints reach back to column 0, so nothing
can be forgotten. That decomposition is therefore prefix enumeration --- correct,
perfectly parallel, and buying distribution rather than pruning.

\subsection{The certificate problem}

The distributed system considered emits no DRAT or LRAT, so an UNSAT verdict
rests on the encoder and the solver both being correct. Four mitigations reduce
what must be trusted: the lemma has independent ground truth; census mode
validates encoder, DAG semantics, enumeration mode and solver core against the
published counts in one run; satisfiable instances catch the dangerous failure
(an over-constrained encoder returns spurious UNSAT, which no amount of UNSAT
running can detect); and certification is available from CaDiCaL or Kissat with
\texttt{drat-trim} for any instance a single node can finish. The sensible
division of labour is scout with the distributed system, certify with
cube-and-conquer plus a DRAT-emitting solver.


%=====================================================================
\section{The extinction model}
\label{sec:extinction}

The investigation's negative results all share a shape: the methods fail because
order 32 is not structurally distinguished. This section argues that the
distinguishing fact is populational rather than structural, and that it is
quantitative.

\subsection{The four-class taxonomy}

Drakakis, Iorio and Rickard partition Costas arrays into four classes:
\emph{generated} (G), obtained by algebraic constructions whose applicability
depends only on the order; \emph{predictably emergent} (PE), arising from
transformations of generated arrays and guaranteed by additional conditions;
\emph{unpredictably emergent} (UE), arising from such transformations but
requiring direct verification; and \emph{sporadic} (S), everything else.

This sharpens language used loosely earlier in this report. The catalogue
generator of Section~\ref{sec:catalogue} produces $G\cup PE$ and part of $UE$; the
16 arrays of order 29 it cannot reach are \emph{emergent}, not sporadic --- they
were found previously by a semi-empirical technique. Order 29 has \emph{zero}
sporadic equivalence classes, and order 28 is the first order above 5 with none.

\subsection{The published sporadic curve}

Table 1 of the order-28 enumeration gives sporadic equivalence-class counts:

\begin{center}
\begin{tabular}{@{}lrrrrrrrrr@{}}
\toprule
order & 13 & 14 & 17 & 19 & 24 & 25 & 26 & 27 & 28 \\
\midrule
total ECs & 1616 & 2168 & 2294 & 1283 & 25 & 12 & 8 & 29 & 89 \\
sporadic ECs & 1602 & 2163 & 2270 & 1281 & 24 & 5 & 2 & \textbf{1} & \textbf{0} \\
\bottomrule
\end{tabular}
\end{center}

Sporadics dominate the middle orders, collapse through the twenties, and vanish
at 28. The extinction conjecture --- that beyond some $N$ only generated and
predictably emergent arrays survive --- is theirs, argued from the order-28
result, from unsuccessful searches for new emergent arrays, and from the
$O(n^3)$ growth in the number of constraints.

\subsection{A first-moment derivation}

The conjecture can be made quantitative. For lag $d$ the $n-d$ differences
$h_{i+d}-h_i$ must be distinct among roughly $2n-2$ available values, a birthday
cost of about $(n-d)^2/(4n)$ in the log. Summing over $d$,
\[
\log \Pr[\,p \text{ is Costas}\,] \;\approx\; -\frac{1}{4n}\sum_{k=1}^{n-1}k^2
\;\approx\; -\frac{n^2}{12},
\]
against the entropy $\log n! \approx n\log n - n$. The populations cross where
$n\log n - n = n^2/12$, that is $n \approx 12(\log n - 1)$; at $n=29$ the right
side is $28.4$.

\begin{observation}[Quadratic collision cost fits the observed sporadic curve]
\label{obs:extinction}
Fitting $\log S(n) = \log(n!/8) - f(n)$ with $f$ quadratic to the sporadic EC
counts above gives
\[
f(n) = 0.08501\,n^2 + 0.173\,n - 3.657,
\]
whose leading coefficient agrees with the derived $1/12 = 0.08333$ to within
$2\%$.
\end{observation}

\begin{center}
\begin{tabular}{@{}lrrrrrrrr@{}}
\toprule
order & 13 & 14 & 17 & 19 & 24 & 25 & 26 & 27 \\
\midrule
sporadic ECs & 1602 & 2163 & 2270 & 1281 & 24 & 5 & 2 & 1 \\
model & 1830 & 2171 & 1941 & 1032 & 25.5 & 8.3 & 2.4 & 0.60 \\
\bottomrule
\end{tabular}
\end{center}

The fit holds across four orders of magnitude, is essentially exact at $n=14$
(2171 against 2163), and places the extinction point --- where the model first
falls below one equivalence class --- at $n=27$, one order before the observed
extinction at 28.

\subsection{Consequences for order 32}

Extrapolating: $0.025$ sporadic ECs at order 29, $0.0042$ at 30, $0.0006$ at 31,
and
\[
\boxed{\;\approx 10^{-4}\ \text{sporadic equivalence classes at order }32.\;}
\]

The contribution from the recognised direct Welch and Golomb construction
families is zero under the taxonomy and applicability conditions used here.
Orders 31 and 32 admit no direct construction: $W_1$ needs $n+1$ prime, $G_2$ needs $n+2$ a prime power, and for
$n=31,32$ none of $32,33,34,35,36$ supplies one in the required position
(Proposition~\ref{prop:polydual}). Order 31 in particular has \emph{no} direct
construction --- $W_1$ would require $p=32$ prime --- yet it has arrays, and rung
$k=1$ identifies their origin exactly --- augmentation of order-30 cores, 8
arrays in one orbit. For order 32, several corresponding extension routes have
been exhausted: rung $k=2$ over all 664 \emph{generated} order-30 cores --- order
30 has never been enumerated, so this clause is catalogue-relative --- together
with the census-complete rungs $k=3$ and $k=4$ over every order-29 and every
order-28 Costas array respectively, the latter at $115{,}087{,}822{,}400$
placements. These close the specified augmentation and deletion-minor channels
but do not control an order-32 array possessing none of those minors.

So the expected order-32 population is $10^{-4}$ from the sporadic term plus zero
from the algebraic term. This explains, without recourse to any structural
obstruction, why no amount of computation has produced an order-32 array: the
methods of Sections~\ref{sec:structural}--\ref{sec:coords} failed not because
they were too weak but because the fact they were searching for is populational.

\begin{remark}[A falsifiable prediction]
\label{rem:falsifiable}
The model predicts that an order-30 enumeration will find \emph{only} generated
and emergent arrays --- $W_1(31)$ and $G_2(32)$ derivatives --- and zero
sporadic equivalence classes. A single sporadic array at order 30 would refute
the extinction model and reopen order 32. This raises the value of F4 from data
collection to hypothesis test.
\end{remark}

\begin{remark}[Honest error bar on the extrapolation]
\label{rem:errorbar}
The point estimates above understate the uncertainty. The fit has residual
standard deviation $0.360$ in log space on 5 degrees of freedom, and the
extrapolation leverage at $n=32$ is $6.62$ against $0.46$--$0.50$ inside the data
range $13\le n\le27$. Proper 95\% prediction intervals are therefore wide:
$[6.1\times10^{-6},\,1.0\times10^{-3}]$ at $n=32$, and for the aggregate
$\sum_{n\ge29}\mathbb E[S(n)]$ a point estimate of $0.030$ with a band of roughly
$[0.007,\,0.14]$. Read as a Poisson probability that no sporadic array ever
appears again, that is $97\%$ as a point estimate but $87\%$--$99\%$ as an
interval. Leave-one-out refits are stable (order-32 predictions from
$7.6\times10^{-5}$ to $1.6\times10^{-4}$), but leave-one-out tests interpolation
at leverage $0.5$ and says nothing about extrapolation at leverage $6.6$.
\end{remark}

\begin{remark}[Extinction versus completeness]
The natural theorem here is not $S_n=0$ for large $n$. ``Sporadic'' denotes what
present knowledge fails to explain, so that statement is contingent on the state
of the literature rather than on mathematics. The well-formed version is a
\emph{completeness} theorem: every Costas array of sufficiently large order
arises from an explicit list of constructions. That is the analogue of the
classification of finite simple groups, and no structural theory of the sporadic
class exists to support it.
\end{remark}

\subsection{Model-form adequacy, and where the forecast's weight actually sits}
\label{sec:adequacy}

The interval in Remark~\ref{rem:errorbar} accounts for estimation uncertainty
\emph{conditional on the quadratic model}. It says nothing about the model form.
With eight observations that distinction is decisive.

\paragraph{Stress test.} Fitting a cubic to the same data:

\begin{center}
\begin{tabular}{@{}lrrrrr@{}}
\toprule
model & RSS & resid.\ sd & $S(30)$ & $S(32)$ & $\sum_{n\ge29}$ \\
\midrule
quadratic & 0.649 & 0.360 & $4.2\times10^{-3}$ & $7.8\times10^{-5}$ & $0.030$ \\
cubic & \textbf{0.256} & \textbf{0.253} & $4.4\times10^{-2}$ & $8.8\times10^{-3}$ & $\sim10^{191}$ \\
\bottomrule
\end{tabular}
\end{center}

\noindent This is a model-adequacy stress test, not a competing forecast. The
cubic \emph{fits better} in sample and extrapolates to eventual explosion rather
than extinction; the two curves agree on the calibration range and diverge by
some 192 orders of magnitude outside it. The cubic is not rejected by goodness of
fit. It is rejected only relative to the proposed collision-cost mechanism, under
which $f(n)=-\log\Pr[\pi_n \text{ Costas}]$ should eventually increase and grow at
least quadratically.

\paragraph{That exclusion is a hypothesis, not a theorem.} It is tempting to
argue that adding a row adds constraints and removes none, so $f$ must be
increasing and convex. The argument fails: $\pi_n\sim\mathrm{Unif}(S_n)$ and
$\pi_{n+1}\sim\mathrm{Unif}(S_{n+1})$ are different probability spaces and are not
nested, deletion with compression does not preserve uniformity, and a superlinear
count of formal constraints does not imply convexity of their aggregate
logarithmic cost, since dependence among them also grows.

\paragraph{What the data say about the hypothesis.} Computing
$f(n)=\log n!-\log C(n)$ at every enumerated order:

\begin{center}
\begin{tabular}{@{}lr@{}}
\toprule
$\Delta f>0$ (monotone) & holds at 27 of 28 consecutive orders \\
$\Delta^2 f>0$ (convex) & holds at 20 of 27, failing at $n=4,8,10,22,25,26,27$ \\
\bottomrule
\end{tabular}
\end{center}

Monotonicity is well supported. \emph{Convexity is not}, and three of its seven
failures are the last three orders before the extrapolation begins. The
construction-free sporadic series shows the same thing: the slope of
$g(n)=\log(n!/8)-\log S(n)$ runs $+2.34,\,+2.76,\,+3.20,\,+3.88,\,+4.79$ and then
\emph{falls} to $+4.17$ and $+3.99$. Those terminal points carry counts of 5, 2
and 1, so Poisson noise of 45--100\% could account for it entirely; but the data
do not support convexity where it matters most.

\paragraph{Convexity is exactly what carries the conclusion.} Dropping it and
continuing $g$ linearly at its terminal slope $4.317$:

\begin{center}
\begin{tabular}{@{}lrrr@{}}
\toprule
model class for $g$ & decreasing phase & infinite horizon & $\Pr[\text{none}]$ on that phase \\
\midrule
superlinear (quadratic $f$) & $0.030$ & $0.030$ & $0.97$ \\
linear at terminal slope & $0.618$ through $n=73$ & \emph{diverges} & $0.54$ \\
concave (cubic fit) & model-dependent & \emph{diverges} & $\approx0$ \\
\bottomrule
\end{tabular}
\end{center}

\noindent The two descriptions must be kept apart. Under the linear model the
step ratio is $(n+1)e^{-s}$, which exceeds 1 once $n\ge74$ at $s=4.317$; the
figure $0.618$ is the expected count accumulated over the \emph{decreasing} phase
$28\le n\le73$ only, and $e^{-0.618}=0.539$ is the corresponding probability on
that finite window. Over the infinite horizon the series diverges, so a Poisson
reading gives probability zero of no further sporadic array. Extinction requires
$g$ to grow \emph{superlinearly}, and nothing in the enumeration establishes that:
under linear growth the sporadic population declines through
the thirties, bottoms out near $n=74$, and resumes growing.

\begin{observation}[Where the forecast's weight sits]
\label{obs:weight}
The extinction forecast rests on two separable components: a structural
modelling assumption, that the effective logarithmic constraint cost is
eventually increasing and grows at least quadratically; and an empirical
calibration, which estimates the rate of that growth over the enumerated range.
The data support the calibration only after the model class has been chosen; they
do not establish the class. Within the class, $\Pr[\text{no further sporadic}]$ is
$0.97$ with an estimation band of $0.87$--$0.99$. Across plausible classes it
ranges from $0.97$ to below $0.54$.
\end{observation}

\paragraph{The constant is not pinned either.} Three levels should be
distinguished: the naive independent-uniform birthday calculation gives
$1/12=0.0833$; correcting for the approximately triangular difference law, whose
effective support is 47.2 slots at order 32 against 62 if uniform, raises it to
about $0.109$; the fitted value is $0.08501$. The gap between the second and third
shows that dependence or finite-order effects materially reduce the effective
collision cost relative to the refined independent model. It is not evidence for
any particular mechanism doing so. The derivation supports the quadratic scale and
order of magnitude far more convincingly than it identifies the constant.

\begin{remark}[The asymptotic theorem the argument needs]
\label{rem:openproblem}
Extinction in the first-moment picture requires the entropy to be dominated:
$\mathbb E[S(n)]=(n!/8)e^{-f(n)}\to0$ summably, that is
\[
f(n)-\log(n!/8)\;\longrightarrow\;+\infty
\]
fast enough, with $f(n)\ge cn^2$ or $f(n)\ge(1+\varepsilon)n\log n$ as concrete
sufficient forms. Eventual monotonicity and eventual convexity are informative
intermediate results but neither implies extinction: a convex function may be
asymptotically linear, and linear $f$ is dominated by the entropy $n\log n$. A
theorem supplying entropy domination would be significant in its own right, is
the missing ingredient in every probabilistic statement above, and --- unlike the
order-32 existence question --- is uniform in $n$, so
Observation~\ref{obs:uniformity} does not apply to it.
\end{remark}

The model is a first-moment heuristic, not a theorem. It says nothing about
whether a particular order-32 array exists, only what the expected number is
under a random model calibrated on nine measured orders. The quadratic scale and order of magnitude were motivated before fitting. The
fitted coefficient lies within $2\%$ of the naive independent-uniform value
$1/12$, but not of the triangular-support refinement near $0.109$; the agreement
should therefore not be read as identifying the asymptotic constant.



%=====================================================================
\section{Transpose-symmetric arrays and a closed branch}
\label{sec:symmetric}

A restriction to symmetric arrays is tractable at order 32 in a way the general
problem is not, and one branch of it is now closed.

\subsection{Transpose symmetry is the only case}

The non-trivial stabilisers in $D_4$ reduce to a single conjugacy class. A
horizontal flip fixes $h$ only if $h_i=h_{n-1-i}$, not injective; a vertical flip
needs $2h_i=n-1$ for all $i$; a $90^\circ$ rotation squares to a $180^\circ$ one;
and the anti-diagonal reflection is conjugate to the main diagonal. The
$180^\circ$ case is the interesting one:

\begin{proposition}
\label{prop:no180}
No Costas array of order $n\ge3$ has $180^\circ$ rotational symmetry.
\end{proposition}
\begin{proof}
If $h_i=n-1-h_{n-1-i}$, the pair $(i,j)$ and its image $(n{-}1{-}j,\,n{-}1{-}i)$
have identical displacement $(j-i,\ h_j-h_i)$. They are distinct unless
$i+j=n-1$, and such pairs exist for $n\ge3$.
\end{proof}

So a search over involutions $h=h^{-1}$ is exhaustive for the symmetric question.
Measured, with three independent implementations agreeing on entered nodes:

\begin{center}
\begin{tabular}{@{}rrr@{}}
\toprule
$n$ & entered nodes & arrays \\
\midrule
16 & 160{,}797 & 40 \\
18 & 864{,}774 & 20 \\
20 & 4{,}710{,}666 & 8 \\
22 & 25{,}853{,}570 & 10 \\
24 & 142{,}822{,}399 & \textbf{0} \\
\bottomrule
\end{tabular}
\end{center}

\noindent The growth ratio is about $5.5$ per step of two orders, projecting about
$1.4\times10^{11}$ nodes at order 32, about one to two core-days on a single core. (This
corrects an earlier figure of $6\times10^{14}$, which applied the ratio once per order.)

\subsection{Fixed points form a Golomb ruler}

\begin{proposition}
\label{prop:fpgolomb}
The fixed points of a transpose-symmetric Costas array form a Golomb ruler.
Hence $\binom f2\le n-1$, and at $n=32$, since the minimum length of an
eight-mark Golomb ruler is 34 and the board offers span 31, $f\in\{0,2,4,6\}$.
\end{proposition}
\begin{proof}
If $h(u)=u$ and $h(v)=v$ with $v-u=d$, the pair has displacement $(d,d)$, so two
fixed-point pairs at equal separation repeat a displacement.
\end{proof}

Verified on all 158 arrays at orders 12--22, with no exceptions; at $n=16$ the
bound correctly predicts the observed absence of $f=6$, which the counting bound
alone permits.

\begin{observation}[True but non-binding]
\label{obs:fpnonbinding}
Implemented as a search prune, the cap fires \emph{zero} times: node counts are
identical with and without it at $n=16,18,20,22$. The incremental Costas check
already subsumes it, since a repeated fixed-point distance \emph{is} a repeated
displacement. The deepest fixed-point count the search reaches is 4 at $n=16$
and 5 at $n=20,22$, against caps of 4 and 6. Recorded so that it is not
implemented twice.
\end{observation}

\subsection{The branch $f=2$, $M_{LL}=0$ is closed}

Let $M=A^++A^-$ record displacements by magnitude, split at 15. The marginals and
transpose symmetry give the identity $X_r-Z_r=-r^2+(2n-1)r-\binom n2$, verified
single-valued on every array at orders 12--22; at $n=32$, $r=15$ it reads
$X-Z=224$, so the small-small quadrant cannot vanish. Exhaustive computation
gives $|\operatorname{Aut}(G_{15})|=2$, identity and reversal, both of which
repeat displacements, so the off-diagonal quadrants cannot vanish either. Only
$M_{LL}=0$ survives, and it forces all transposition lengths and all
fixed-point separations below 16, whence $f\in\{2,4\}$ and, at $f=2$, the
transposition lengths are exactly $1,\dots,15$.

Writing $m_d=a_d+b_d$ for the doubled centre of the length-$d$ transposition,
every cross-cycle displacement has diagonal coordinates
$(m_e-m_d,\ \pm(e-d))$ or $(m_e-m_d,\ \pm(e+d))$. Two sieves follow:
\textbf{FG}, that $m_{d+k}-m_d$ is injective in $d$ for every gap $k$; and
\textbf{FS}, that $m_e-m_d$ is injective among pairs with $d+e=s$.

\begin{theorem}
\label{thm:f2closed}
No transpose-symmetric Costas array of order 32 has $f=2$ and $M_{LL}=0$.
\end{theorem}

The exhaustive classification is
\[
32{,}769{,}242 \ \xrightarrow{\ \mathrm{FG}\ }\ 9{,}180
\ \xrightarrow{\ \mathrm{FS}\ }\ 4
\ \xrightarrow{\ \text{Costas}\ }\ 0 ,
\]
over all 168 admissible fixed-point pairs. FG alone removes $99.972\%$. The four
survivors form two reflection orbits under $\rho(i)=31-i$, which acts by
$(p,q)\mapsto(31-q,31-p)$ and $m_d\mapsto62-m_d$, and each orbit carries an
explicit one-line collision certificate.

\emph{Orbit $\mathcal O_C$}, representative $p=10$, $q=18$: here
$m_2-m_9=m_{15}-m_8=-5$ while $2-9=-(15-8)=-7$, so the segments $(0,1)=(a_9,a_2)$
and $(21,22)=(b_8,b_{15})$ both have displacement $(1,-6)$. This is a
\emph{sign-reversed} gap relation, $e-d=-(e'-d')$, which FG as stated does not
test --- one segment joins lower endpoints and the other upper endpoints.

\emph{Orbit $\mathcal O_A$}, representative $p=12$, $q=18$: here
$2q-m_{13}=m_2-m_{11}=-11$ with $-13=-(11+2)$, so $(2,3)=(a_{11},b_2)$ and
$(17,18)=(a_{13},q)$ both have displacement $(1,-12)$. This one involves a fixed
point and is not internal to the transposition cycles.

Both certificates were verified independently by reconstructing the involutions
and comparing directed displacement vectors.

\begin{remark}
The result is a finite computer-assisted classification, not a deductive proof
that FG and FS admit only four centre sequences. What it does supply is an
exclusion that no longer rests on a terminal ``zero Costas'' count: each
reflection class fails by a displayed identity. The first-failure histogram over
the fixed-sum fibres peaks at $s=11$, not at the maximal-cardinality fibres
$s=15,16,17$, so fibre size does not predict strength --- though this is a
first-failure count under a fixed testing order, not an order-independent measure.
\end{remark}


\subsection{The branch $f=4$, $M_{LL}=0$ is empty}
\label{sec:f4closed}

\begin{theorem}
\label{thm:f4closed}
No transpose-symmetric Costas array of order 32 has $f=4$ and $M_{LL}=0$.
\end{theorem}

Under $M_{LL}=0$ the four fixed points form a Golomb ruler of diameter at most
15, the 14 transpositions have distinct lengths in $\{1,\dots,15\}$ with exactly
one omitted length $r$, and the parity of the fixed-point set is coupled to $r$
by $a\equiv r+1\pmod2$. The census over that parameter space is

\[
N_F=2{,}735 \;\longrightarrow\; N_{F,r}=20{,}360
\;\longrightarrow\; N_{\mathrm{domain}}=20{,}360
\;\longrightarrow\; N_{\mathrm{cover}}=\mathbf 0 ,
\]

where $N_F$ counts reflection-normalised four-mark Golomb rulers of diameter at
most 15, $N_{F,r}$ the parity-compatible ruler/omitted-length pairs,
$N_{\mathrm{domain}}$ those whose every required transposition length retains a
non-empty centre domain, and $N_{\mathrm{cover}}$ those admitting an
endpoint-disjoint cover. No candidate reaches the fixed-gap, signed fixed-gap,
fixed-sum or anchor tests, let alone the Costas oracle.

\begin{observation}[Local viability, global incompatibility]
\label{obs:f4local}
The algebraic filters are \emph{non-binding} here. The anchor-centre exclusion
set has $|\mathcal A(F)|=10$ for \emph{every} ruler --- no coincidences anywhere
--- and yet $N_{\mathrm{domain}}=N_{F,r}$ exactly: not one parameter pair loses a
domain. Total domain sizes run from 89 to 196 with median 143, so these are not
starved instances. The parity condition merely halves the range of $r$ and leaves
a flat survivor distribution. Non-existence appears only when the local choices
must be made simultaneously, at the endpoint exact-cover layer.
\end{observation}

It is tempting to summarise this as ``four marks crowded into a window of
diameter 15 cannot be complemented by fourteen distinct-length edges''. That is a
useful reading but not a deductive mechanism; what is proved is the exhaustive
zero-cover census under the domains specified above.

\begin{corollary}
\label{cor:mll0}
No transpose-symmetric Costas array of order 32 satisfies $M_{LL}=0$.
\end{corollary}
\begin{proof}
Under $M_{LL}=0$, $f\in\{0,2,4,6\}$. The case $f=0$ needs 16 transpositions with
distinct lengths drawn from $\{1,\dots,15\}$; $f=6$ needs a six-mark Golomb ruler
inside a window of diameter 15, against a minimum length of 17; $f=2$ is
Theorem~\ref{thm:f2closed}; $f=4$ is Theorem~\ref{thm:f4closed}.
\end{proof}

\subsection{The $Z=1$ layer is empty}
\label{sec:z1cells}

Corollary~\ref{cor:mll0} moves the residual symmetric problem to the large-large
mass $Z$: every remaining candidate has $Z\ge1$. Since $M$ is symmetric,
$Z\equiv D_L\pmod 2$, so $Z=1$ splits into one long transposition (branch T) or
one long fixed pair (branch F), each with its admissible fixed-point counts.
All seven resulting labels are now closed.

\begin{center}
\begin{tabular}{@{}llll@{}}
\toprule
branch & $f$ & certificate & method \\
\midrule
$Z{=}1\mathrm T$ & 0 & $36\to1\to0$ & one residue, killed by Costas \\
$Z{=}1\mathrm F$ & 2 & $36\to0$ & structural \\
$Z{=}1\mathrm T$ & 2 & $72\to0$ & full Costas \\
$Z{=}1\mathrm T$ & 6 & --- & deductive: OGR(6)$=17>15$ window \\
$Z{=}1\mathrm F$ & 6 & $728\to0$ & structural \emph{and} full Costas \\
$Z{=}1\mathrm F$ & 4 & $19{,}405\to0$ & structural \emph{and} full Costas \\
$Z{=}1\mathrm T$ & 4 & $2{,}735\to0$ & full Costas \\
\bottomrule
\end{tabular}
\end{center}

\begin{theorem}
\label{thm:z1empty}
No transpose-symmetric Costas array of order 32 has $Z=1$. With
Corollary~\ref{cor:mll0}, every remaining transpose-symmetric order-32 candidate
satisfies $Z\ge2$.
\end{theorem}

A parity relation governs each cell: writing $a$ for the number of even fixed
points and $c$ for the number of odd-length transpositions, counting
cross-parity edges gives $c\equiv a\pmod2$, whence $o(R)\equiv a+L$ in branch T
and $o(R)\equiv a$ in branch F. With $R=\varnothing$ this forces $L$ even; with
one omitted length the omission absorbs $L$'s parity; with two omissions in
$Z{=}1\mathrm T$, $f=4$ it reads $L\equiv a+r+s\pmod2$.

The two largest censuses were $Z{=}1\mathrm F$, $f=4$ at $5.89\times10^{10}$
structural nodes over 19{,}405 roots, and $Z{=}1\mathrm T$, $f=4$ at
$2.49\times10^{10}$ nodes over 2{,}735 reflection-canonical ruler roots. For the
latter, the parameter space --- 2{,}735 rulers, 297{,}838 long placements and
15{,}642{,}886 parity-admissible $(F,L,u,r,s)$ records --- was derived
independently \emph{before} the solver existed, and the campaign's internal
traversal totals reproduce those figures exactly, so the census covered the
intended premise space rather than merely producing one record per outer root.

\begin{remark}[Positive-arm validation]
\label{rem:t48}
A non-existence census can fail silently if its terminal layer is unreachable.
The $Z{=}1\mathrm T$, $f=4$ solver was therefore tested in the positive
direction, under an explicitly weakened specification: removing \emph{only} the
incremental $Z_p+z_{\mathrm{add}}>1$ prune and the terminal $Z_p\ne1$ rejection,
and leaving the endpoint cover, centre sum, length spectrum and recursion
untouched. The unchanged recursion then reaches accepted terminal states,
increments its cover counter and emits completions: on root F0000, 3{,}498{,}614
survivor lines during a 60-second bounded run --- not an exhaustive count and
not a claimed cardinality --- of which 2{,}000 sampled objects were
independently revalidated and all 2{,}000 passed, with the first 200{,}000
emitted records all distinct.

This establishes that the recursive terminal, counter and emission paths operate
when completions exist. It does \emph{not} establish the correctness of the
enforced $Z=1$ budget, which is covered separately by comparing the incremental
counter against direct recomputation. Naming the specification under which the
path was exercised matters: an earlier version of the same solver shipped a
self-test that never passed, because its fixture was rejected before reaching
the assertion it was written to make.
\end{remark}

\begin{remark}[Two withdrawn implementations]
\label{rem:withdrawn}
A fast full-Costas prototype reported a $3{,}600$--$7{,}900\times$ speedup on the
expensive Cell-3 roots. External audit found two independent defects. The first
loses collisions among displacements created by a single placement: the
implementation checks all new displacements against the global state and only
then inserts them, so two equal new displacements both pass and collapse to one
bit. It was caught by the population invariant
$\texttt{displacements}=\binom{\text{points}}{2}$ failing by exactly one per
transposition placed. The second declares the rollback snapshot \texttt{static},
so every recursion level shares one buffer and a parent restores its most recent
descendant's state. The first defect under-prunes; the second \emph{over}-prunes
and can therefore manufacture false zeros. The benchmark and its zero counts are
withdrawn in full.

This is recorded because of how it was found. In both cases the author's reading
of the source concluded the code was correct, and in both cases a measurement
disagreed --- the population invariant for the first defect, an independent
reimplementation for the second. It is the eighth such episode in this
investigation, and the pattern is consistent enough to state as a working rule:
here, an argument that code is correct carries no evidential weight against a
test that says otherwise.
\end{remark}



%=====================================================================
\section{Proof inventory and certification status}
\label{sec:certificates}

Thirty-six formal claims appear in this report. This section records the
certification state so that it is visible rather than implicit, and separates
claims by the kind of certificate each would need.

\textbf{Three claims are now machine-checked.} Sixteen theorems typecheck under
Lean 4.15.0 with \textbf{no Mathlib, no \texttt{sorry} and no
\texttt{native\_decide}}; ten depend on no axioms whatever and six on
\texttt{propext} and \texttt{Quot.sound} only, both part of Lean's standard
trusted base. Avoiding Mathlib was deliberate: everything is formulated over core
\texttt{List} and \texttt{Int} rather than \texttt{Finset}, so the development
checks in seconds from a bare toolchain with no dependency resolution.
Reproduce with \texttt{make lean}.

\begin{center}
\begin{tabular}{@{}lll@{}}
\toprule
claim & Lean & content \\
\midrule
A5 telescoping identity & \textbf{proved} & \texttt{telescope}, \texttt{T\_symm}, with \texttt{win\_shift}, \texttt{win\_split} \\
A7 overdetermination $\not\Rightarrow$ inconsistency & \textbf{proved} & \texttt{cut\_labelling\_consistent} \\
A3 prime-power obstruction at 32 & \textbf{proved} & \texttt{no\_prime\_power\_33/34/35/36}, \texttt{prime\_power\_31/32} \\
A2 $T_m$ soundness & core proved & \texttt{tm\_interval\_sound}: disjoint outer intervals $\Rightarrow$ no completion \\
A1 diagonal parity floor & core proved & \texttt{parity\_identity}, \texttt{diagTriples\_240}, \texttt{parity\_floor\_arith} \\
\bottomrule
\end{tabular}
\end{center}

\noindent A2 needs two further pieces to close Lemma~\ref{lem:tmsound}: the
extremal lemma (\emph{$k$ distinct elements of a finite integer set have sum
between its $k$ smallest and its $k$ largest}) and the derivation that each
computed interval is an outer bound. A1 needs the combinatorial half --- that
$2\mid D$ iff the endpoints share value parity, and that the 240 diagonal triples
biject with same-parity index pairs. Both arithmetic cores are done.

\begin{center}
\begin{tabular}{@{}llr@{}}
\toprule
tier & character & count \\
\midrule
A & deductive; formalisable from the statement alone & 13 \\
B & computational; needs \texttt{decide}, \texttt{native\_decide} or a verified checker & 9 \\
C & heuristics, model fits and measurements; not certification targets & 13 \\
\midrule
 & \textbf{proved} & \textbf{3} \\
 & \textbf{partially proved} & \textbf{2} \\
\bottomrule
\end{tabular}
\end{center}

\subsection{Tier A: deductive claims}

Ordinary mathematics; a Lean proof is a matter of effort rather than research.
The thirteen are the diagonal parity floor (Thm~\ref{thm:parityfloor}), soundness
of the $T_m$ rule (Lem~\ref{lem:tmsound}), the automatic loss-free
parameterisation (Prop~\ref{prop:polydual}), offset invariance of interpolation
degree, the telescoping identity and $T_m=T_{n-m}$, the fact that one off-curve
point forces full degree, the $K_s$ counterexample showing overdetermination does
not imply inconsistency, truncation destroying associativity, a Costas array not
being a representation, the corrected board-network characterisation, the 2-adic
tower domination theorem, the $\ell_\infty$ collision bound, and corner deletion
acting by translation.

\subsection{Tier B: computational claims}

True statements about finite objects, established by search. A certificate means
\texttt{decide} (kernel-checked, small instances only), \texttt{native\_decide}
(fast but kernel-bypassing, and worth recording as a distinct status), or a
verified checker consuming an externally produced certificate.

\begin{center}
\begin{tabular}{@{}llll@{}}
\toprule
claim & scale & realistic route & priority \\
\midrule
Order-29 census reconstruction & $10^2$ & \texttt{decide} & \textbf{in progress} \\
The missing orbit & $10^3$ & \texttt{decide} & \\
Order-28 catalogue completeness & $10^3$ & \texttt{decide} & \\
Heredity is corner-driven & $10^4$ & \texttt{native\_decide} & \\
Order-28 census classification & $10^6$ & \texttt{native\_decide} & \\
F3 computational result & $3\times10^5$ & checker; soundness half is Tier A & \\
$m_2\ge1$ & $2^{17}$ cases & checker, or the structural proof & \\
Census-complete rung $k=3$ & $5.7\times10^8$ & certificate $+$ checker & \\
Rung results $k=2$, $k=4$ & $1.2\times10^{11}$ & checker architecture only & \\
\bottomrule
\end{tabular}
\end{center}

The order-29 census reconstruction is the cheapest real certificate available and
underwrites the quantifier in Theorem~\ref{thm:rung3census}. The highest-value
item is Lemma~\ref{lem:tmsound}: certifying it converts the F3 soundness result
from ``no false positives observed over $7{,}632$ completable prefixes'' into
``cannot produce a false positive'', making the empirical half redundant. Its one
genuinely new ingredient is that $k$ distinct elements of a finite integer set
have sum between that of its $k$ smallest and its $k$ largest.

\begin{remark}[What a Tier B certificate would and would not establish]
\label{rem:tierb}
A machine-checked Tier B result certifies that a particular algorithm, run on
particular data, produced a particular answer. It does not certify that the
algorithm implements the intended mathematical object; that link is Tier A work.
The two implementation defects in Section~\ref{sec:tmprop} are the cautionary
case. Both were internally consistent and deterministic, and both would have
passed a certificate of their own output. What exposed them was testing the
positive arm --- checking that no completable configuration is rejected --- not
any form of proof.
\end{remark}

\subsection{Order-29 census: the first Tier B certificate}
\label{sec:census29lean}

This was identified above as the cheapest real certificate available and the one
underwriting the quantifier in Theorem~\ref{thm:rung3census}. It is now largely
constructed, in Lean 4.15.0 with no Mathlib, no \texttt{sorry} and no
\texttt{native\_decide}.

\paragraph{What kernel-checks.} All 23 published equivalence-class
representatives, each as an individual theorem
\texttt{checkSeed29 census29SeedNN = true} discharged \texttt{by decide}. Each
verifies that representative \emph{and all eight of its $D_4$ images} as order-29
Costas permutations: length, range, permutation uniqueness and fixed-lag
difference uniqueness. Total elaboration is about 230\,s, roughly 10\,s per seed.
The two orbit-level theorems also compile individually: the expansion has
cardinality exactly 164 (10\,s), and every member of it is a Costas permutation
(120\,s).

\paragraph{Three findings about the mechanism.} First, per-seed chunking is what
makes kernel decision feasible at all. A single aggregate over the 23
representatives re-runs every seed inside one reduction and exceeds the
elaborator's budget; it is also redundant, since the per-seed theorems are the
certificate of record. Second, the binding constraint is \emph{heartbeats}, not
recursion depth: the orbit theorems fail at the default 200{,}000-heartbeat limit
regardless of \texttt{maxRecDepth}, and pass with
\texttt{set\_option maxHeartbeats 4000000}. Third, kernel reduction is
memory-hungry --- the complete module exhausts a 4\,GB container even when every
constituent theorem compiles alone, so the development is split into a seed
module and an orbit module. The orbit module needs a larger machine than was
available here; nothing in it is unproven, only unbuilt in one sitting.

\paragraph{The boundary, which matters more than the build.} This certifies the
supplied representatives and their 164-element orbit expansion. It does
\emph{not} prove that the catalogue is exhaustive. Exhaustiveness of the order-29
census is an external premise inherited from the published enumeration, and it is
precisely the premise on which Theorem~\ref{thm:rung3census}'s universal quantifier
rests. So the certificate closes the verification half --- every array claimed to
be Costas is one --- while the completeness half remains a citation. This is the
same distinction drawn in Remark~\ref{rem:tierb}, and the paper should state it
in one line rather than let the presence of a Lean certificate imply more than it
delivers.

\paragraph{An earlier attempt, and why it failed.} A first version placed the
whole computation --- generation, $D_4$ closure, deduplication and Costas
checking over all 164 arrays --- under a single \texttt{by decide}. It returns
\texttt{true} instantly under \texttt{\#eval} and fails in the kernel, because
\texttt{decide} does not use the interpreter: it asks the kernel to reduce the
computation by \texttt{whnf}, with no compiled code and no fast paths through
\texttt{List}. Roughly $10^6$ comparisons is well past what that will do.
\texttt{native\_decide} would finish at once and is exactly the kernel-bypassing
route the choice of \texttt{decide} was meant to avoid. Agreement under
\texttt{\#eval} is the positive-arm test, not the certificate.



%=====================================================================
\section{External review and residual directions}
\label{sec:review}

Five independent assessments of this report were commissioned in French, German,
Russian, Japanese and Chinese. Two observations about them are worth recording
before their content.

\textbf{No new literature and no breakthroughs.} None of the five cites a paper,
result, author or dataset absent from the bibliography here. All reason from this
report's own record. Five independent reviews turning up no external development
is itself weak evidence that the literature is well covered.

\textbf{One stale status, correctly caught.} Item F6b was marked blocked on data
after the order-29 census had already been reconstructed. The screen has now been
run and the item closed (Observation~\ref{obs:rw0}).

The reviews converge strongly on priorities already ranked here --- publication,
an incremental propagator, certifying the order-29 census first, order-30 as a
hypothesis test --- which is corroboration rather than news. Several describe the
$T_m$ soundness lemma and the dyadic near-Costas test as open; both are complete,
so those sections were read from an earlier snapshot. What follows are the five
directions that were \emph{not} already in this programme.

\paragraph{F17. $RW_0$ as an operation class, not an array class. (Weeks.)}
The two $RW_0$ classes were analysed as arrays; the \emph{operation} producing
them was not. Appending a blank row to an exponential $W_1$ array, cyclically
shifting, and adjoining a column is the one concrete Costas-producing move known
to lie outside corner surgery --- and the entire augmentation ladder of
Section~\ref{sec:ladder} is corner surgery. The questions are whether the
operation generalises to orders 30--32, whether its output admits a local Costas
test, and whether it closes under composition with corner operations.
\emph{Continue if:} the operation extends to order 30 and produces arrays the
generated catalogue misses. This is the only item here that attacks the residual
class directly rather than closing another channel.

\paragraph{F18. A residual-class specification with an order-sensitivity gate.
(Months.)} Define the residual class by exclusion --- no order-28 minor, no
order-29 minor, outside the recognised construction families, outside the
generated order-30 augmentation cores --- then search for invariants on it,
calibrated against the complete order-28 and order-29 censuses. The success
criterion is the novel part: not that an invariant classifies known arrays
cleanly, but that it \emph{treats orders 31, 32 and 36 differently}. That turns
Observation~\ref{obs:uniformity} from a caution into a design specification.
\emph{Abandon if:} no candidate invariant separates 31 from 32 on the censuses
available.

\paragraph{F19. A residual-only SAT encoding. (Months.)}
Distinct from the generic SAT encoding closed in Section~\ref{sec:sat}: build an
encoding that assumes the exclusions already proved --- no order-28 deletion
minor, no order-29 deletion minor --- and carries $T_m$, fixed-lag, parity and
endpoint bounds as native propagators, then splits by cube-and-conquer with
DRAT/LRAT certificates. This uses the theorems of Section~\ref{sec:ladder} as
search structure rather than starting from an unconstrained instance.
\emph{Abandon if:} median shard time or tail risk is unmanageable at order 24--26
pilots.

\paragraph{F20. A hybrid portfolio solver. (Months.)}
Only monolithic propagation configurations have been benchmarked. The F3$'$
measurements are themselves the argument against that: incremental-only runs
$53$--$63\times$ faster in wall time while exploring $5.5$--$8.1\times$ more
nodes (Section~\ref{sec:f3prime}), so the shallow and deep regimes plainly want
different treatment. Cheap incremental bitmask at shallow depth, $L_3$ or
$L_3+T_m$ at mid depth by predicted benefit, SAT for highly constrained
subproblems, with fully reproducible decision logs. \emph{Gate:} identical
solution sets at orders 11--13 and wall time better than every single strategy.

\paragraph{F21. Pre-registered competing extinction models with an ex-ante
order-30 prediction. (Weeks.)} Section~\ref{sec:extinction} acknowledges model
uncertainty; this goes further. Register in advance the $n^2$ first-moment model,
the linear-terminal-slope model, a cubic, and a model stratifying
construction-derived from sporadic classes; perform leave-one-\emph{order}-out
predictive checks; and publish a prior interval for the order-30 sporadic count
\emph{before} the computation runs. That converts order 30 from a test of one
model into a discriminating comparison among several, and directly addresses the
weakness identified in Observation~\ref{obs:weight}.

\paragraph{Honourable mentions.} Two further suggestions are good but concern
dissemination rather than the object. A methodology paper on establishing finite
negative computational results --- the four disciplines of
Section~\ref{sec:discipline} together with the two propagator defects of
Section~\ref{sec:tmprop} --- would reach a considerably wider audience in SAT,
constraint programming and exhaustive search than the Costas literature. And a
stratified sampling campaign before any full order-30 enumeration, maximising the
probability of finding a sporadic counterexample per unit of compute by targeting
structurally likely regions first, would be a rational ordering of that budget.

%=====================================================================
\section{Synthesis: what is now known about order 32}

\begin{quote}
The investigation does not resolve existence at order 32. It proves two
census-complete deletion-minor exclusions, classifies the generated augmentation
channel from order 30, reconstructs and independently verifies the order-29
census, and introduces a sound telescoping-sum propagator that passes a
prospective complete-population benchmark with substantial marginal coverage at
low additional measured work. The remaining existence problem consists of arrays
outside the tested construction and deletion-minor channels, and no current
structural classification controls that residual class.
\end{quote}

\label{sec:synthesis}

\begin{enumerate}[label=\textbf{S\arabic*.},leftmargin=*]

\item \textbf{The periodic branch is partly closed.} No circular Costas array of
order 32 exists, because 33 is composite (cited). Single periodicity remains open
pending the Golomb--Taylor conjecture.

\item \textbf{The two informative augmentation rungs are closed.} No order-32
Costas array is a two-point augmentation of any algebraic order-30 core
(Theorem~\ref{thm:rungs}); the minimum defect over all such
configurations is 2, attained in exactly one dihedral orbit.

\item \textbf{The score-2 object has a complete genealogy.} It is the unique
minimum-cost one-point extension of the unique derived order-31 Costas orbit
(Observation~\ref{obs:genealogy}), which is itself the unique one-dot
augmentation of the algebraic order-30 cores. This is a far sharper description
than ``Welch-30 plus two dots'' and explains the rigidity observed from every
direction: the modulo-16 template has one defect-2 lift with a gap to 13; the
augmentation census has one defect-2 orbit with defects 1 and 3 empty; bounded
Hamming surgery finds no repair.

\item \textbf{The published sporadic count first reaches zero at order 28 and
remains zero at order 29.} A first-moment model with approximately quadratic
constraint cost tracks the observed decline
(Observation~\ref{obs:extinction}) and predicts about $10^{-4}$ sporadic
equivalence classes at order 32, while the algebraic term at that order is empty.
This is a \emph{conditional population forecast, not a proof of permanent
extinction}: the data calibrate the rate within the chosen model class but do not
establish the required superlinear asymptotic growth
(Observation~\ref{obs:weight}). Within that caveat it is the first account that
explains the failure of every method in this report rather than adding to it ---
the distinguishing fact about order 32 is populational, not structural.

\item \textbf{Rung $k=4$ is closed against the complete order-28 census.}
89 units, all \texttt{EXHAUSTED}, $115{,}087{,}822{,}400$ core placements, zero
Costas arrays, independently merged and verified. Because order 28 has no
sporadic arrays this is the first rung quantifying over every core of its order:
\textbf{two independent census-complete augmentation channels are now closed}:
no order-32 Costas array has an order-29 Costas deletion minor
(Theorem~\ref{thm:rung3census}), and none has an order-28 Costas deletion minor.
Neither result implies the other without a hereditary-minor theorem, which
is not established and is not supported by the order-28 measurements
(Remark~\ref{rem:independent}).

\item \textbf{One unconditional bound exists.} Every order-32 permutation has at
least 112 diagonal defects divisible by two (Theorem~\ref{thm:parityfloor}). The
associated mod-4 minimisation does not bear on existence
(Observation~\ref{obs:m2}).

\item \textbf{Costas arrays are almost never hereditary}
(Theorem~\ref{thm:heredity}), and essentially all inherited sub-arrays come from
corner deletion. This caps the reach of the entire constructed-and-augmented
branch.

\item \textbf{The order-28 census is completely classified by coordinate
system} (Theorem~\ref{thm:census28}): 672 arrays are Welch-direct, 32 are Golomb
descendants visible only in dual-exponential coordinates over $\F_{31}$ and
$\F_{2^5}$, and a single orbit is a corner augmentation and generic in both. The
governing asymmetry is that corner deletion preserves the algebraic coordinates
while corner addition destroys the parent's degree-one parameterisation in the
augmented orbit studied here.

\item \textbf{The prime-power shadow has an automatic-parameterisation dual}
(Proposition~\ref{prop:polydual}). For the standard Welch or Golomb
parameterisation to fill an order-$N$ board automatically, without deletion or
window selection, one needs $q-1=N$ or $q-2=N$. At $N=32$ this requires field
orders 33 or 34, neither of which is a prime power. Over the next available
field, $q=37$, the random-output window factor is
$32!/36^{32}=4.1542\times10^{-15}$; this is a sampling heuristic, not a
non-existence theorem for specially structured maps.

\item \textbf{Every order-insensitive structural relaxation tested fails to
distinguish order 32, and the common cause
is identified} (Observation~\ref{obs:uniformity}). Relation algebra, matroid
theory, affine and toric invariants, the 2-adic tower and finite-field coordinate
families were each implemented and measured, not merely argued against.

\item \textbf{Telescoping-sum propagation supplies a substantial new
finite-depth obstruction.} On the complete order-12 six-leading-column benchmark,
$L_3+T_m$ detects $117{,}880$ of $296{,}656$ infeasible prefixes, including
$40{,}806$ beyond the frozen baseline, for about $5\%$ additional declared
propagation work (Theorem~\ref{thm:f3}). It passes the prospectively frozen gate
and generalises across the additional complete benchmarks tested. It has
\emph{not} been shown to reduce the asymptotic search exponent or full-solver
runtime at higher orders.

\item \textbf{The exhaustive cost is unchanged} at roughly $10^3$--$10^4$ modern
core-years for order 32, with a measured growth ratio of $4.80$ per order.

\item \textbf{Therefore any order-32 Costas array would lie outside the
currently recognised construction families and outside the tested extension
closure.} This is the precise form of the statement; ``sporadic'' is
catalogue-relative rather than an intrinsic mathematical class, and a
hypothetical order-32 array could exhibit a construction not yet recognised. The
residual class is defined by absence of known structure and is therefore immune
to algebraic attack by construction. S4 makes it larger than previously supposed,
not smaller.

\end{enumerate}

%=====================================================================
\section{Future work, ranked}

For the transpose-symmetric programme specifically, the validation gates are
closed --- including the positive-arm test of Remark~\ref{rem:t48} --- and one
assurance objective remains: reproduction of the $Z=1$ closures by a genuinely
independent second implementation. No further testing of the existing solvers
can supply that, since implementation independence is the property at issue. The
mathematical frontier is $Z\ge2$.

\label{sec:future}

Ranked by expected value per unit effort. Each entry states the deliverable, the
cost, the expected outcome, and the criterion for abandonment.

\subsection*{Completed since the initial investigation}

\paragraph{F2 --- rung $k=3$ reproducibility. \textbf{closed}.} Re-run over all
148 arrays of the corrected algebraic generator: exhausted, zero Costas arrays,
minimum defect 4.

\paragraph{F1b --- order-29 census acquisition. \textbf{closed}.} The complete
census was reconstructed from Table 2 of the enumeration paper and validated
independently, reproducing all published figures: 164 arrays, 23 equivalence
classes, 5 symmetric, per-method counts $4/64/24/24/32/16$.

\paragraph{Census-complete rung $k=3$. \textbf{closed}.} $565{,}836{,}800$
placements, zero Costas arrays, minimum defect 4, yielding
Theorem~\ref{thm:rung3census}.

\subsection*{Tier 1 --- do these}

\paragraph{F1. Publish the present result set. (Highest value; weeks.)}
The results above are novel, verified and at risk of evaporating. A single paper
covers: the heredity measurement with the corner-dot identification
(Theorem~\ref{thm:heredity}) --- cheap, novel and explanatory, since it accounts
for why every construction above order 24 proceeds by corner adjunction; the
augmentation-ladder classification with the order-31 uniqueness and the genealogy
of $N_{32}$; the catalogue method of Section~\ref{sec:catalogue}, which
reproduces the order-28 census exactly and lands on precisely $164-16=148$ at
order 29, constituting an independent verification of the published enumerations;
and the negative methodology of Section~\ref{sec:structural}, which is the
component most likely to be useful to others because it saves them a year.
\emph{Abandon if:} never.

\paragraph{F3$''$. Make $L_3$ cheap enough to pay for itself. (Weeks.)}
F3$'$ is complete (Section~\ref{sec:f3prime}): $T_m$ passes the scientific gate
and fails the operational one, and so does $L_3$ itself --- incremental-only
search is $53$--$63\times$ faster in wall time despite exploring up to $8.1\times$
more nodes. The bottleneck is singleton shaving, not $T_m$, which adds only $4\%$
work. The experiment worth running is whether an incremental, watched-support
$L_3$ with a sparse invocation schedule can beat the bare incremental check at
all; only if it can does the $26$--$32\%$ tree reduction from $T_m$ become
operationally available. The operational target must be
\emph{incremental-only}, not $L_3$: a stack that beats $L_3$ while remaining
$50\times$ slower than the production solver is of no use for enumeration. The
candidate levers, in order of expected return, are incremental state updates
replacing repeated matching and shaving rebuilds; caching propagation results for
repeated domain states; sparser $T_m$ invocation behind a cheap interval-width
precheck; and running $L_3$ at selected depths only with incremental checks
between. Any scheduling tuning must use development orders disjoint from the
evaluation orders, unlike the a-priori schedule used in F3$'$.
\emph{Abandon if:} no $L_3$ variant reaches wall-time parity with
incremental-only at orders 13--15.

\paragraph{F4. Order-30 exhaustive enumeration --- now a hypothesis test. (300--500 modern core-years.)}
Never attempted. It would produce the first complete census above order 29 and
would settle empirically whether the sporadic population is extinct at 30 or
merely thinning --- the single most informative fact bearing on S8. The
infrastructure (fail-closed work units, deterministic prefix manifests, merge
with coverage verification) is largely built. Cost estimate from the measured
growth law anchored on the order-29 figure of 366.55 CPU-years. \emph{Abandon
if:} a GPU or FPGA implementation cannot reach within an order of magnitude of
the required throughput after a bounded engineering effort.

\paragraph{F5. Solver throughput engineering. (Weeks; prerequisite for F4.)}
The measured single-core rate is $4.34\times10^6$ nodes/s at order 32. Published
FPGA work reports speedups of about $40\times$ over a GPU and $4.44\times$ over
the fastest reported software implementation. A GPU implementation of the bitmask
DFS with full dihedral symmetry reduction is the realistic route. Deliverables:
deterministic prefix generation, one file per work unit, completion records,
result hashes, a manifest of expected prefixes, and a merge tool that rejects
missing or duplicated work. \emph{Abandon if:} throughput does not improve by at
least $20\times$ over the single-core baseline.

\paragraph{F6b. Screen the order-29 emergent arrays. \textbf{closed}.}
Formerly blocked on data; the census reconstruction removed the blockage and the
screen has been run. Both $RW_0$ classes are coordinate-generic
(Observation~\ref{obs:rw0}).

\paragraph{F7. Attack the Golomb--Taylor conjecture. (Open since 1984.)}
The only prime-sensitive mechanism in the literature. Proving that a singly
periodic Costas sequence must be Welch immediately gives no singly periodic
order-32 Costas array, since Welch requires $p-1=32$ and 33 is composite. Moreno
proved that the stronger shifting property characterises Welch; the task is to
close the gap between single periodicity and the shifting property. \emph{Abandon
if:} nothing new is found in six months; it is a recognised open conjecture and
partial progress is publishable.

\paragraph{F8. Formalise the local-versus-global separation as a methodological
result. (Weeks.)}
The data already exists: non-extendable partial Costas configurations at every
order (minimum inclusion-maximal sizes $3,4,4,5,5$ for $n=4,\dots,8$), the
measured detection rates of RA and vertex-domain propagation, and the explicit
instance refuted by RA and not by L3. The publishable statement is that pointed
relation-algebraic extension does not certify global thinness for Sidon-type
conditions, quantified. This generalises well beyond Costas.

\paragraph{F9. Welch wrap-parity compression of the completed census.
(Days.)}
The wrap correction for a $W_1(31)$ core is $(30\epsilon_x, 31\epsilon_y)$ with
$\vtwo(30)=1$ and $\vtwo(31)=0$, so a horizontal wrap changes the column
difference by $\pm30\equiv2\pmod4$ while a vertical wrap flips the parity of the
height difference. This gives four parity buckets and is a legitimate way to
compress the already-completed $k=2$ enumeration into a human-readable covering
proof. It is a presentation device, not a discovery mechanism: the enumeration is
done. \emph{Abandon if:} the covering argument does not close within a week.

\paragraph{F10. Measure algebraic proof complexity. (Days.)}
``Prove algebraically'' has a precise meaning: a Nullstellensatz or
Positivstellensatz certificate for infeasibility, with certificate degree as the
difficulty parameter. Test instances already exist --- the non-extendable partial
configurations such as $\{(1,4),(2,6),(3,3),(5,1),(6,5)\}$ in the $8\times8$
grid. Compute the minimum refuting degree at orders 8 through 12 and observe the
growth. \emph{Abandon if:} the degree grows linearly by order 12, which would
show algebraic proof complexity is the wrong tool.

\subsection*{Tier 4 --- do not pursue}

\paragraph{F11. Deepening the ladder to $k\ge5$.} Removed by
Heuristic~\ref{cor:ladder-value}: these rungs exclude classes the heredity
measurement says are empty \emph{a priori}. The row-side DFS of
Section~\ref{sec:antiramsey}, which would make them feasible, is therefore not
worth implementing for this purpose.

\paragraph{F12. Any further uniform partition or modular quotient.} Closed by
Observation~\ref{obs:uniformity} and quantified throughout
Section~\ref{sec:partition}: no static palette, magnitude, residue, parity or
valuation partition binds, because every horizontal class is at most half full.

\paragraph{F13. Distributed SAT as an accelerator.} Closed by the calibration in
Section~\ref{sec:sat}. Retain only as a scout and a cross-implementation check,
and only with the certificate caveat understood.

\paragraph{F14. Unstructured inverted searches over low-degree permutation
polynomials.} Closed as a discovery strategy, but not by a universal
non-existence theorem. No automatically loss-free Welch or Golomb
parameterisation exists at order 32, because the required field orders 33 and 34
are not prime powers (Proposition~\ref{prop:polydual}). At the next available
field, $q=37$, a uniform-output model assigns one prescribed width-32 injective
window the probability $32!/36^{32}=4.1542\times10^{-15}$, so uninformed sampling
cannot resolve the relevant regime. Independently, the complete order-12 census
shows the non-linear degree body to be statistically indistinguishable from a
matched random control ($\chi^2=1.454$, $p=0.483$), the $d=7$ excess to be
insignificant ($p=0.215$), and the degree-1 class alone to be cleanly separated
($p=4.7\times10^{-4}$). Together these close unstructured low-degree sampling as
a practical programme. They do \emph{not} exclude a future structured family
equipped with an exact mechanism forcing its image into the bounded window. A
reported sweep of roughly 250 candidates returning zero hits is consistent with
any hypothesis and establishes nothing.

\paragraph{F14a. Conditioned dyadic lifting via survivor hypergraphs. (Days;
the one dyadic thread still open.)} The all-escape form is closed by
Observation~\ref{obs:escape}: it is inconsistent for Costas arrays and random
permutations alike, so $\Lambda_k$ carries no signal. But the corrected statistic
does separate (Observation~\ref{obs:muk}): Costas arrays lift at about the fourth
percentile of balanced re-lifts while random permutations lift at the median. The
decisive follow-up is cheap and well defined --- measure the same percentile for
near-Costas permutations of defect 1 and 2. \emph{Continue if:} near-Costas
arrays also lift low, which would make the statistic predictive rather than a
restatement of $S_6=\varnothing$, and would give a search heuristic worth
encoding. \emph{Abandon if:} only exact Costas arrays lift low, since the signal
is then tautological.

\paragraph{F14b. Exact determination of $m_2$.} Closed as an order-32
programme by Observation~\ref{obs:m2}: $Z_2^{\mathrm{diag}}$ does not separate
Costas arrays from random permutations, and minimising it moves away from the
region real Costas arrays occupy. The decision sequence --- a zero-survivor
feasibility model over 240 five-bit forbidden-pattern constraints, then
successive bounds $\sum_\omega z_\omega\le B$ with exact survivor indicators, with
$q_0=0$ as the one safe normalisation --- is correctly specified and remains a
clean finite combinatorial problem in its own right. It should not be carried
under the order-32 banner.

\paragraph{F15. Additive-node polynomial screening.} Closed by calibration: the
statistic assigns known algebraic arrays the same maximal degree, full support and
zero rank deficiency as random permutations, and is therefore incapable of
detecting structure of any kind.

\paragraph{F16. Sharpening the density theorem to reach order 32.} The density
result currently available does not reach an integer non-existence bound at order
32. A stronger theorem would need sufficient order-sensitive numerical strength;
it cannot be ruled out merely because arrays exist at neighbouring orders
(Observation~\ref{obs:uniformity}). Original note follows. Closed by
Observation~\ref{obs:uniformity}: any smooth density bound giving $c(32)=0$ also
gives $c(31)=0$, which is false.

%=====================================================================
\section{Software artefacts}
\label{sec:artefacts}

\begin{longtable}{@{}l p{0.62\textwidth}@{}}
\toprule
Artefact & Description \\
\midrule
\endhead
\texttt{costas\_fast.cpp} & Bitmask exact DFS, one \texttt{uint64\_t} per
  distance; completion status; symmetry note for prefix mode. \\
\texttt{costas\_rung.cpp} & Augmentation-rung engine. Modes \texttt{selftest}
  (reproduces rungs $k=1,2$), \texttt{cat} (generate and census-check a
  catalogue), \texttt{run} (sweep a work unit). Core-first pruning, bucketed
  pair order, generation-stamp collision detection, canonical-orbit mode. \\
\texttt{merge\_rung.py} & Fail-closed merge: refuses an exhaustive verdict unless
  all units report \texttt{EXHAUSTED}, agree on parameters and catalogue hash,
  and cover $0..N{-}1$ exactly once; independently re-scores every reported
  configuration. \\
\texttt{cat28.txt} & The census-complete order-28 catalogue, 712 arrays in 89
  dihedral orbits. \\
\texttt{ra\_closure.cpp} & Relation-algebra closure on 63-bit edge masks with
  Régin GAC per distance class; modes for pointed sweeps, ground-truth detection
  and the score-2 prefix diagnostic. \\
\texttt{prop.py} & Vertex-domain propagator, levels L1--L3. \\
\texttt{gf32\_expanded.cpp} & Expanded $\F_{32}$ affine-power sweep with
  independent coordinate orderings, matched random baseline and $D_4$
  deduplication. \\
\texttt{dyadic\_profile.py} & Verifies the diagonal parity floor identically
  against random permutations, computes the $Z_1\ldots Z_6$ filtration profile of
  any permutation, and calibrates $Z_2^{\mathrm{diag}}$ for a census against a
  matched random control. \\
\texttt{f3\_complete.cpp} & The F3 protocol run: RA path-consistency and $L_3$ baselines, the $T_m$-coupled candidate, frozen gate, deterministic cost metric and soundness check. \\
\texttt{tm\_propagator.cpp} & Earlier F3 propagator suite: $L_1$/$L_2$/$L_3$
  baselines over position--value domains and the $T_m$-coupled propagator,
  benchmarked over complete infeasibility populations with soundness checked on
  every feasible configuration. \\
\texttt{gen\_benchmark12.cpp} & Regenerates the order-$n$ infeasibility
  population by exhaustive completion search. \\
\texttt{escape\_system.py} & Builds the survivor hypergraph and the affine
  escape system $H_kb_k=c_k$ at each dyadic level, reports rank, consistency,
  solution-space dimension and the balanced-lift count $\Lambda_k$, and
  calibrates a census against a matched random control. \\
\texttt{costas\_algebraic\_screen.py} & Single-exponential screen, pure Python.
  Fail-closed validation, $D_4$ canonicalisation, analytic null, and a
  calibration self-test asserting that Welch arrays recover degree 1 while
  printing the additive-node statistic alongside for comparison. \\
\texttt{dual\_exponential\_screen.py} & Dual-exponential screen with GF($p$) and
  GF($2^m$) arithmetic, so Golomb families over extension fields are visible.
  Classifies a catalogue into single-exponential, dual-exponential and residual
  classes and cross-references provenance against generated families. \\
\texttt{gen\_costas\_dagster.py} & CNF and DAG generator: \texttt{lemma},
  \texttt{neighbourhood}, \texttt{census}, \texttt{selftest}. \\
\texttt{check\_dagster\_models.py} & Independent model decoder and verifier,
  including census model counting. \\
\bottomrule
\end{longtable}

%=====================================================================
\begin{thebibliography}{99}

\bibitem{drakakis27} K.~Drakakis, S.~Rickard, J.~K.~Beard, R.~Caballero,
F.~Iorio, G.~O'Brien and J.~Walsh, \emph{Results of the enumeration of Costas
arrays of order 27}, IEEE Trans.\ Inform.\ Theory \textbf{54} (2008).

\bibitem{drakakis29} K.~Drakakis, F.~Iorio, S.~Rickard and J.~Walsh,
\emph{Results of the enumeration of Costas arrays of order 29}, IEEE Trans.\
Inform.\ Theory \textbf{57} (2011).

\bibitem{golombtaylor} S.~W.~Golomb and H.~Taylor, \emph{Constructions and
properties of Costas arrays}, Proc.\ IEEE \textbf{72} (1984).

\bibitem{golombmoreno} S.~W.~Golomb and O.~Moreno, \emph{On periodicity
properties of Costas arrays and a conjecture on permutation polynomials}, IEEE
Trans.\ Inform.\ Theory \textbf{42} (1996).

\bibitem{moreno92} O.~Moreno, \emph{Survey of results on signal patterns for
locating one or multiple targets}, and the shifting-property characterisation
(1992).

\bibitem{warnke} L.~Warnke, B.~Correll and C.~Swanson, \emph{The density of
Costas arrays decays exponentially}, IEEE Trans.\ Inform.\ Theory \textbf{69}
(2023), 575--581.

\bibitem{ddc} S.~R.~Blackburn, T.~Etzion, K.~M.~Martin and M.~B.~Paterson,
\emph{Two-dimensional patterns with distinct differences --- constructions,
bounds, and maximal anticodes}, IEEE Trans.\ Inform.\ Theory \textbf{56} (2010),
1216--1229.

\bibitem{ddc2} S.~R.~Blackburn, T.~Etzion, K.~M.~Martin and M.~B.~Paterson,
\emph{Distinct difference configurations: multihop paths and key predistribution
in sensor networks}, IEEE Trans.\ Inform.\ Theory \textbf{56} (2010),
3961--3972.

\bibitem{dsn} J.~A.~Dias da Silva and M.~B.~Nathanson, \emph{Maximal Sidon sets
and matroids}, arXiv:math/0504226 (2005).

\bibitem{maddux} R.~D.~Maddux, \emph{Relation algebras of every dimension},
J.~Symbolic Logic \textbf{57} (1992); and \emph{Relation Algebras}, Elsevier
(2006).

\bibitem{tarskigivant} A.~Tarski and S.~Givant, \emph{A Formalization of Set
Theory without Variables}, AMS Colloquium Publications \textbf{41} (1987).

\bibitem{hh} R.~Hirsch and I.~Hodkinson, \emph{Relation Algebras by Games},
North-Holland (2002); and \emph{Representability is not decidable for finite
relation algebras}, Trans.\ AMS \textbf{353} (2001).

\bibitem{kowalski} T.~Kowalski, \emph{Representability of Ramsey relation
algebras}, Algebra Universalis (2015).

\bibitem{almmaddux} J.~F.~Alm, R.~D.~Maddux and J.~Manske, \emph{Chromatic
graphs, Ramsey numbers and the flexible atom conjecture}, Electron.\ J.\
Combin.\ (2008).

\bibitem{lamleung} T.~Y.~Lam and K.~H.~Leung, \emph{On vanishing sums of roots
of unity}, J.\ Algebra \textbf{224} (2000).

\bibitem{ptn} M.~J.~H.~Heule, O.~Kullmann and V.~W.~Marek, \emph{Solving and
verifying the Boolean Pythagorean triples problem via cube-and-conquer}, SAT
(2016).

\bibitem{schur5} M.~J.~H.~Heule, \emph{Schur number five}, AAAI (2018).

\bibitem{csplib} CSPLib Problem 076: Costas arrays.

\bibitem{oeis} OEIS A008404 (number of Costas arrays of order $n$) and A001441
(inequivalent Costas arrays under the dihedral group).

\end{thebibliography}

\end{document}
